% !TEX root = ../main.tex
\chapter{Principal Component Analysis and Factor Models}\label{ch:pca}
\begin{paracol}{2}

\begin{sources}
\begin{tabularx}{\linewidth}{@{}l l L@{}}
\toprule
\textbf{Lecture} & \textbf{Speaker} & \textbf{Content}\\
\midrule
2024 L9 & S.~Andreev (Two Sigma) & PCA as a rotation of coordinates; toy examples, robustness, PCA versus OLS; demeaning and variance normalisation; eigenvalue diagnostics, in- and out-of-sample variance explained; PCA on financial time series; US Treasury curve (level, slope, curvature, regime indicators); PC-neutral curves and flies, leverage.\\
2013 L15 & P.~Kempthorne & Linear factor models as cross-sectional and time-series regressions; Sharpe's single-index model, macroeconomic and fundamental (BARRA, Fama--French) models; factor-mimicking portfolios; statistical factor models: factor analysis (maximum likelihood, EM, likelihood-ratio test) and PCA; case study on Treasury yields.\\
\bottomrule
\end{tabularx}

\smallskip
\textbf{Files.} Slides \file{mit18_642_f24_lec09.pdf} (the European-indices example, slides 26--34, was
not covered in class); Jupyter notebook \file{feddataanalysis.ipynb} with the data file
\file{fedbondyielddata_20230911.csv} (distributed in 2024 week 5); readings
\file{mit18_642_f24_pset_analysis.pdf} (P.~Kempthorne, \emph{Principal Components Analysis in Finance:
Modeling the Dynamics of Interest Rates}, 2016) and \file{mit18_642_f24_pset_rcode.pdf} (its R code and a
5-year PCA of Treasury yields); 2013 slides \file{MIT18_S096F13_lecnote15.pdf} and case study
\file{MIT18_S096F13_CaseStudy7.pdf}.
\textbf{Problem sets.} 2024 PS2 Problem~3 and PS3 Problem~4; 2013 PS7 (all problems).
\end{sources}

\section*{Motivation}
Financial data come in large and highly redundant blocks: the yields of thirty Treasury maturities,
the returns of five hundred stocks, a dozen European equity indices. When the 10-year yield rises the
5- and 30-year yields almost always rise too, so thirty numbers per day carry perhaps three numbers'
worth of information. This chapter develops two complementary ways of exploiting the redundancy:
\begin{itemize}
  \item \textbf{Principal component analysis (PCA)}: a rotation of coordinates such that the first few
        axes carry most of the variance. It is model-free (unsupervised learning, no dependent
        variable), pure linear algebra, and fast \ts{24}{9}{3:05}.
  \item \textbf{Factor models}: the statistical model $\bm x_t=\bm\alpha+B\bm f_t+\bm\epsilon_t$ with a few
        common factors and uncorrelated asset-specific noise. It reduces the $m(m+1)/2$ entries of a
        covariance matrix to $O(mK)$ parameters, and is the backbone of risk models and of the
        CAPM/APT view of returns \ts{13}{15}{0:00}.
\end{itemize}
The 2024 lecture is the practitioner's view (a fixed-income quant at Two Sigma: workflow, pitfalls, and
how PCA turns into tradeable portfolios); the 2013 lecture is the statistician's view (models, estimation,
tests). Both analyse daily changes of US Treasury yields and find the same three factors: \emph{level},
\emph{slope} and \emph{curvature}. The mathematics needed is the spectral theorem, the SVD and
Perron--Frobenius of \cref{ch:linalg} (\cref{thm:spectral,thm:svd,thm:pf}) and least squares
(\cref{ch:regression}).

\begin{coursenote}[Notation across the sources]
The 2024 slides use a data matrix $X$ of size $N\times p$ (rows = observations) and loadings $w_k$, $W$;
the 2013 slides use $m$ assets, $T$ periods, a data matrix of size $m\times T$ (columns = periods) and
eigenvectors $\gamma_k$, $\Gamma$; the 2016 reading uses a tenors$\times$days matrix $A$. Here:
$p$ variables and $N$ observations with $X\in\R^{N\times p}$ in the PCA sections, $m$ assets and $T$ periods in
the factor-model sections, and $\gamma_k$, $\Gamma$ for eigenvectors (loadings) throughout.
\end{coursenote}

% ------------------------------------------------------------------
\section{Principal components of a random vector}\label{ch07:sec-pop}

Let $\bm x=(x_1,\dots,x_p)\T$ be a random vector with $\E[\bm x]=\bm\mu$ and $\Cov[\bm x]=\Sigma$. For any
weight vector $\bm w$, $\Var(\bm w\T\bm x)=\bm w\T\Sigma\bm w\ge0$: $\Sigma$ is symmetric positive semi-definite
(\cref{sec:symmetric}; 2024 PS2 Problem~3). If $\bm x$ are asset returns and $\bm w$ portfolio weights,
$\bm w\T\Sigma\bm w$ is the portfolio variance. By the spectral theorem (\cref{thm:spectral})
\[
  \Sigma=\Gamma\Lambda\Gamma\T=\sum_{k=1}^p\lambda_k\,\bm\gamma_k\bm\gamma_k\T,\qquad
  \lambda_1\ge\lambda_2\ge\dots\ge\lambda_p\ge0,\qquad \Gamma\T\Gamma=I_p .
\]

\subsection{Variance maximisation}

\begin{definition}[Principal components]\label{ch07:def-pca}
The \emph{first principal direction} is a unit vector $\bm w_1$ maximising $\Var(\bm w\T\bm x)=\bm w\T\Sigma\bm w$
subject to $\|\bm w\|=1$. The $k$-th principal direction $\bm w_k$ maximises the same variance subject to
$\|\bm w\|=1$ and $\bm w\perp\bm w_1,\dots,\bm w_{k-1}$ \ts{13}{15}{1:07:08}. The \emph{principal component
variables} (scores, PC factors) are $p_k=\bm w_k\T(\bm x-\bm\mu)$; the entries of $\bm w_k$ are the
\emph{loadings} of PC~$k$, and $\lambda_k/\sum_j\lambda_j$ is the \emph{fraction of variance explained}.
\end{definition}

\begin{theorem}[PCA is the spectral decomposition of $\Sigma$]\label{ch07:thm-pca}
With the notation above:
\begin{enumerate}[(i)]
  \item the $k$-th principal direction can be taken to be the eigenvector $\bm\gamma_k$, and
        $\Var(p_k)=\lambda_k$ (if the eigenvalues are distinct, $\bm w_k=\pm\bm\gamma_k$);
  \item $\bm p=\Gamma\T(\bm x-\bm\mu)$ has $\E[\bm p]=\bm 0$ and $\Cov[\bm p]=\Lambda$: the PCs are uncorrelated;
  \item $\bm x=\bm\mu+\Gamma\bm p=\bm\mu+\sum_{k=1}^p p_k\,\bm\gamma_k$;
  \item total variance is preserved: $\sum_i\Var(x_i)=\tr\Sigma=\tr\Lambda=\sum_k\lambda_k=\sum_k\Var(p_k)$.
\end{enumerate}
\end{theorem}
\begin{proof}
(i) Expand $\bm w=\sum_jc_j\bm\gamma_j$; then $\|\bm w\|^2=\sum_jc_j^2=1$ and $\bm w\T\Sigma\bm w=\sum_j\lambda_jc_j^2\le\lambda_1$,
with equality for $\bm c=\bm e_1$ (and only for $\bm c=\pm\bm e_1$ if $\lambda_1>\lambda_2$). Inductively, if
$\bm w_j=\bm\gamma_j$ for $j<k$, the constraint $\bm w\perp\bm\gamma_1,\dots,\bm\gamma_{k-1}$ forces $c_1=\dots=c_{k-1}=0$, so
$\bm w\T\Sigma\bm w=\sum_{j\ge k}\lambda_jc_j^2\le\lambda_k$, attained at $\bm\gamma_k$.

(ii) $\E[\bm p]=\Gamma\T(\E[\bm x]-\bm\mu)=\bm 0$ and $\Cov[\bm p]=\Gamma\T\Sigma\Gamma=\Gamma\T\Gamma\Lambda\Gamma\T\Gamma=\Lambda$
\ts{13}{15}{1:01:42}.

(iii) Multiply $\bm p=\Gamma\T(\bm x-\bm\mu)$ by $\Gamma=(\Gamma\T)^{-1}$.

(iv) $\tr(\Gamma\Lambda\Gamma\T)=\tr(\Lambda\Gamma\T\Gamma)=\tr\Lambda$ \ts{13}{15}{1:09:11}.
\end{proof}

Equivalently, with a Lagrange multiplier, $\nabla_{\bm w}[\bm w\T\Sigma\bm w-\lambda(\bm w\T\bm w-1)]=0$ gives
$\Sigma\bm w=\lambda\bm w$: the stationary points of the Rayleigh quotient $\bm w\T\Sigma\bm w/\bm w\T\bm w$ are the
eigenvectors and the stationary values the eigenvalues (\cref{sec:symmetric}).

\begin{intuition}[A rigid rotation of the data cloud]
$\bm x\mapsto\Gamma\T(\bm x-\bm\mu)$ moves the origin to the centre of mass and rotates the axes; distances are
unchanged \ts{13}{15}{1:02:46}. Andreev's toy example: a one-dimensional data set embedded in the plane
along the $x$-axis, the diagonal or the $y$-axis looks two-dimensional, but PCA immediately returns one
direction with all the variance and a second with none \ts{24}{9}{5:11}; add a little noise in the
perpendicular direction and $\lambda_2$ becomes small but non-zero \ts{24}{9}{7:22}. PCA returns two things:
\emph{directions} (eigenvectors) and their \emph{importance} (eigenvalues). For a physicist, $\Sigma$ is the
second-moment tensor of the cloud of points (unit masses); the inertia tensor $I=\tr(\Sigma)\,\mathbb 1-\Sigma$ has
the same eigenvectors, so the principal directions are the \emph{principal axes of inertia}, PC1 being the
axis with the smallest moment of inertia (the direction along which the cloud is most elongated).
\end{intuition}

\begin{remark}[Uniqueness]\label{ch07:rem-sign}
If $\bm\gamma_k$ is an eigenvector, so is $-\bm\gamma_k$: the sign of each PC is arbitrary (the 2024 slides:
``sign does not matter''), and different software may return different signs for the same data. If
$\lambda_k=\lambda_{k+1}$, any orthonormal basis of the eigenspace is equally valid: the directions are not
determined at all. This is the mathematical root of the robustness problems of \cref{ch07:sec-robust}.
\end{remark}

\begin{example}[Two variables with equal variances]\label{ch07:ex-2d}
For $\Sigma=\sigma^2\bigl(\begin{smallmatrix}1&\rho\\ \rho&1\end{smallmatrix}\bigr)$ with $\rho>0$: $\lambda_{1,2}=\sigma^2(1\pm\rho)$ with
$\bm\gamma_1=\frac1{\sqrt2}(1,1)\T$ (the sum) and $\bm\gamma_2=\frac1{\sqrt2}(1,-1)\T$ (the difference), whatever $\rho$ is;
PC1 explains the fraction $(1+\rho)/2$. As $\rho\to0$ the eigenvalues merge: an isotropic ``blob'' has no
preferred direction and moving a single data point can swing the estimated PC1 by any angle
\ts{24}{9}{12:44}. (The general bivariate case is 2013 PS7 Problem~1.)
\end{example}

\subsection{Geometry: best subspace, and PCA versus regression}\label{ch07:sec-geom}

\begin{proposition}[Best $K$-dimensional approximation]\label{ch07:prop-best}
Let $V\in\R^{p\times K}$ have orthonormal columns, so that $VV\T$ is the orthogonal projection onto their span.
Then
\[
  \E\bigl\|(\bm x-\bm\mu)-VV\T(\bm x-\bm\mu)\bigr\|^2=\tr\Sigma-\tr(V\T\Sigma V)\ \ge\ \sum_{j>K}\lambda_j ,
\]
with equality for $V=\Gamma_K=[\bm\gamma_1\cdots\bm\gamma_K]$. The first $K$ PCs span the $K$-dimensional subspace that
minimises the mean squared \emph{perpendicular} distance of the data from it.
\end{proposition}
\begin{proof}
With $\bm y=\bm x-\bm\mu$, $\|(I-VV\T)\bm y\|^2=\bm y\T(I-VV\T)\bm y$ since $I-VV\T$ is a projection; taking expectations gives
$\tr\Sigma-\tr(V\T\Sigma V)$. Now $\tr(V\T\Sigma V)=\sum_j\lambda_jc_j$ with $c_j=\|V\T\bm\gamma_j\|^2$. Each
$c_j\in[0,1]$ (a projection does not increase length) and $\sum_jc_j=\tr(V\T\Gamma\Gamma\T V)=\tr(V\T V)=K$. Hence
\[
  \sum_j\lambda_jc_j-\sum_{j\le K}\lambda_j=\sum_{j\le K}\lambda_j(c_j-1)+\sum_{j>K}\lambda_jc_j
  \le\lambda_K\Bigl(\sum_{j\le K}(c_j-1)+\sum_{j>K}c_j\Bigr)=0,
\]
and $V=\Gamma_K$ gives $c_j=1$ for $j\le K$, $0$ otherwise.
\end{proof}
For a sample this is the Eckart--Young statement of \cref{sec:svd}: the truncated SVD of the centred data
matrix is its best rank-$K$ approximation.

\paragraph{PCA is not regression} \ts{24}{9}{8:28}. OLS of $y$ on $x$ minimises \emph{vertical} squared
distances, OLS of $x$ on $y$ \emph{horizontal} ones, PCA \emph{perpendicular} ones (\cref{ch07:fig-ols}).
For standardised variables with correlation $\rho$ the three lines through the centre have slopes $\rho$,
$1/\rho$ and $1$ \ts{24}{9}{9:31}. OLS presupposes a dependent variable (a causal or predictive
direction); PCA treats all coordinates symmetrically and says nothing about causality. Use regression to
predict $y$ from $x$, PCA to describe the joint correlation structure \ts{24}{9}{10:34}. Two further
differences: PCA is equivariant under rotations of the data (rotate the cloud, the PCs rotate with it)
\ts{24}{9}{11:38}, but \emph{not} under rescaling of individual variables, whereas OLS predictions are
unaffected by changes of units. Hence the covariance-versus-correlation question of
\cref{ch07:sec-corr}.

\begin{figure}[ht]
\centering
\begin{tikzpicture}[scale=1.05,>=Stealth]
  \draw[->,black!50] (-2.9,0) -- (2.9,0) node[right,black]{\small $x$};
  \draw[->,black!50] (0,-2.9) -- (0,2.9) node[above,black]{\small $y$};
  \draw[black!30,dashed,rotate=45] (0,0) ellipse ({2*sqrt(1.6)} and {2*sqrt(0.4)});
  \foreach \p in {(0.24,0.54),(-0.06,-0.74),(-0.26,-0.95),(0.31,1.55),(-0.30,-0.63),(0.78,0.90),(0.36,-0.54),
    (0.21,0.89),(-1.24,-1.01),(-1.85,-2.14),(-1.79,-1.12),(-1.16,-0.28),(0.41,0.19),(-2.53,-1.83),(0.19,0.34),
    (-1.44,-1.15),(-0.84,-1.11),(1.41,0.18),(0.21,1.07),(-0.40,-0.21),(0.36,0.39),(-1.11,-0.44),(1.74,-0.32),
    (1.19,0.92),(1.08,-0.37),(0.32,0.85),(0.03,0.78),(0.17,0.84),(1.83,0.54),(0.46,-0.04),(0.38,-0.76),
    (-0.40,-0.29),(1.23,1.90),(-1.22,-1.31),(0.95,-1.18),(-0.27,-0.12),(1.63,1.70),(-0.12,-0.29),(-0.03,1.52)}
    \fill[black!45] \p circle (1.2pt);
  \draw[thick,thmcol] (-2.8,-1.68) -- (2.8,1.68) node[right,font=\scriptsize]{OLS $y\sim x$ (slope $\rho$)};
  \draw[thick,defcol] (-1.68,-2.8) -- (1.68,2.8) node[above,font=\scriptsize]{OLS $x\sim y$ (slope $1/\rho$)};
  \draw[very thick,intcol] (-2.55,-2.55) -- (2.55,2.55) node[right,font=\scriptsize]{PC1 (slope 1)};
  % residuals of one point
  \fill[black] (-0.47,1.72) circle (1.8pt) node[above left,font=\scriptsize]{$P$};
  \draw[thick,thmcol] (-0.47,1.72) -- (-0.47,-0.282);
  \draw[thick,defcol] (-0.47,1.72) -- (1.032,1.72);
  \draw[thick,intcol] (-0.47,1.72) -- (0.625,0.625);
\end{tikzpicture}
\caption{Forty standardised points with sample correlation $\rho=0.6$ and the $2\sigma$ ellipse. For the point
$P$ the three segments are the distances minimised by OLS of $y$ on $x$ (vertical), OLS of $x$ on $y$
(horizontal) and PCA (perpendicular); cf.\ slide 8 of \file{lec09} \ts{24}{9}{9:31}.}\label{ch07:fig-ols}
\end{figure}

\subsection{Positive correlations and the market factor}\label{ch07:sec-pf}
If all covariances are positive, $\Sigma_{ij}>0$ (typical for yields of different maturities or for stocks in
one market), the Perron--Frobenius theorem (\cref{thm:pf}) applies to $\Sigma$: $\lambda_1$ is simple and
$\bm\gamma_1$ can be chosen with \emph{all entries positive} (2024 PS2 Problem~3). PC1 is then a long-only
portfolio: the ``market'' or ``level'' factor, and all other PCs, being orthogonal to a positive vector,
have loadings of both signs: they are long/short spreads.

\begin{example}[Equicorrelated assets (not discussed in class)]\label{ch07:ex-equi}
For $p$ assets with common variance $\sigma^2$ and common correlation $\rho>0$,
$\Sigma=\sigma^2[(1-\rho)I+\rho\,\bm 1\bm 1\T]$. Then $\lambda_1=\sigma^2(1+(p-1)\rho)$ with $\bm\gamma_1=\bm1/\sqrt p$ (the
equal-weighted portfolio), and $\lambda_2=\dots=\lambda_p=\sigma^2(1-\rho)$. PC1 explains
$(1+(p-1)\rho)/p\to\rho$ as $p\to\infty$: even a modest average correlation produces a dominant first
component in a large universe, while the remaining $p-1$ eigenvalues are degenerate, so the individual
higher PCs are not identified (\cref{ch07:rem-sign}).
\end{example}

\begin{exercise}[PS2 (2024), Problem 3: principal components, special case]\label{ch07:ex-ps2}
A portfolio holds $n$ assets with fixed weights $\bm w=(w_1,\dots,w_n)\T$; over a given horizon the assets have
rates of return $\bm X=(X_1,\dots,X_n)\T$ with mean $\bm\mu$ and covariance $\Sigma=(\sigma_{ij})$, and the portfolio return
is $Y=\bm w\T\bm X$.
\begin{enumerate}[(a)]
  \item Show that $\E[Y]=\bm w\T\bm\mu$ and $\Var[Y]=\bm w\T\Sigma\bm w$.
  \item Prove that $\Sigma$ is positive semi-definite, for any random variables $X_1,\dots,X_n$.
  \item Suppose all entries of $\Sigma$ are positive. Prove that the largest eigenvalue $\lambda_{\max}$ has multiplicity one and
      an eigenvector $\bm v$ that can be chosen with all components positive.
  \item Write the first principal component variable $PC1$ in terms of $\bm X$, $\bm\mu$ and $\bm v$. What are its mean and
      variance?
\end{enumerate}
\end{exercise}

\begin{exercise}[not from the course]\label{ch07:ex-equi-proof}
Prove the claims of \cref{ch07:ex-equi}: for $\Sigma=\sigma^2[(1-\rho)I+\rho\bm1\bm1\T]$ find all eigenvalues and eigenvectors, and
the fraction of variance explained by PC1. For which $\rho$ is $\Sigma$ positive semi-definite? Relate the result to
\cref{thm:pf} when $\rho>0$.
\end{exercise}

% ------------------------------------------------------------------
\section{Sample PCA, SVD and practical choices}\label{ch07:sec-sample}

\subsection{Data matrix and singular value decomposition}
Stack $N$ observations (days) of $p$ variables (tenors, assets) in $X\in\R^{N\times p}$, row $i$ being $\bm x_i\T$.
With $\bar{\bm x}=X\T\bm1/N$, the centred matrix and sample covariance are
\[
  X_c=X-\bm 1\bar{\bm x}\T=\bigl(I_N-\tfrac1N\bm1\bm1\T\bigr)X,\qquad S=\frac{1}{N-1}X_c\T X_c
\]
(the 2013 slides divide by $N$; it does not matter for the eigenvectors) \ts{13}{15}{41:21}.

\begin{proposition}[PCA via the SVD]\label{ch07:prop-svd}
Let $X_c=UD\hat\Gamma\T$ be a thin SVD: $U\in\R^{N\times p}$ with $U\T U=I_p$, $D=\diag(d_1\ge\dots\ge d_p\ge0)$,
$\hat\Gamma\in\R^{p\times p}$ orthogonal. Then:
\begin{enumerate}[(i)]
  \item $S=\hat\Gamma\hat\Lambda\hat\Gamma\T$ with $\hat\lambda_k=d_k^2/(N-1)$: the right singular vectors are the loadings;
  \item the $N\times p$ matrix of PC scores $P=X_c\hat\Gamma=UD$ has zero column means and orthogonal columns;
        column $k$ (the time series of PC~$k$) has sample variance $\hat\lambda_k$;
  \item $X_c=P\hat\Gamma\T=\sum_kd_k\bm u_k\hat{\bm\gamma}_k\T$, and truncation after $K$ terms is the best rank-$K$
        approximation of $X_c$;
  \item $\rank X_c\le\min(N-1,p)$: with $N\le p$ observations at most $N-1$ eigenvalues are non-zero.
\end{enumerate}
\end{proposition}
\begin{proof}
(i) $X_c\T X_c=\hat\Gamma DU\T UD\hat\Gamma\T=\hat\Gamma D^2\hat\Gamma\T$.

(ii) $P=UD\hat\Gamma\T\hat\Gamma=UD$, so
$P\T P=D^2$; $\bm1\T P=\bm1\T X_c\hat\Gamma=\bm0$ because $\bm1\T X_c=\bm 0$.

(iii) is \cref{thm:svd} and the Eckart--Young
theorem.

(iv) $\bm1\T X_c=\bm0$ means the $N$ rows of $X_c$ are linearly dependent.
\end{proof}

Andreev computes the covariance matrix and diagonalises it (``we need the covariance matrix for other
things anyway''); the SVD gives the same result without forming $S$ \ts{24}{9}{18:02}, \ts{13}{15}{1:06:05}.
Numerically the SVD is preferable because forming $X_c\T X_c$ squares the condition number; in cost both
routes are $O(Np^2)$ (plus $O(p^3)$ for the eigen-decomposition).

\begin{coursenote}[Erratum: slides 10--11 of \file{lec09}]
\begin{itemize}
  \item ``$\Sigma=X^TX$ is the covariance matrix'' needs the factor $1/N$ (or $1/(N-1)$); without it the SVD
      relation reads $X\T X=W M\T MW\T$ with $M=\diag(\sqrt{(N-1)\lambda_k})$, not $\diag(\sqrt{\lambda_k})$.
  \item ``Computation of the covariance matrix is $O(N^2p^2)$'' is wrong: $X\T X$ costs $O(Np^2)$, the same order as
      the thin SVD; the advantage of the SVD is accuracy, not asymptotic cost.
  \item ``$W$ is the eigenvectors of
      $X$'': $W$ holds the right singular vectors of $X$, i.e.\ the eigenvectors of $X\T X$.
  \item Slide 11 calls
      $\Sigma=W\Lambda W\T$ the ``SVD decomposition'': it is the spectral decomposition (for a PSD matrix it coincides with the SVD).
\end{itemize}
\end{coursenote}

\begin{exercise}[\file{lecnote15}, slide 32: PCA via the SVD]\label{ch07:ex-lecsvd}
Let $X^*$ be the $m\times T$ row-centred data matrix, $\hat\Sigma_x=\frac1TX^*X^{*\prime}$, and $X^*=VDU'$ its SVD with $V$ ($m\times m$)
orthogonal, $U$ ($T\times m$) with orthonormal columns and $D=\diag(d_1\ge\dots\ge d_m\ge0)$. Show that
$\hat\Lambda=\frac1TD^2$, $\hat\Gamma=V$, and that the matrix of sample principal components is $P=\hat\Gamma'X^*=DU'$.
\end{exercise}

\subsection{A worked example: six days of Treasury yields}\label{ch07:sec-sixdays}
The 2016 reading (\file{pset_analysis}) runs PCA on the smallest possible data set: FRED constant-maturity
yields of 9 Treasury securities (3 months to 20 years) on the first six business days of 2001. Daily
changes in basis points, minus their row means, form the $9\times5$ matrix $A$ of \cref{ch07:tab-sixdays}
(tenors $\times$ days, the transpose of $X_c$, so the roles of $U$ and $\hat\Gamma$ are swapped: loadings are the
\emph{left} singular vectors of $A$).

\begin{table}[ht]
\centering\small
\caption{Six-day example: centred yield changes $A$ (bp) and the PCA of \file{pset_analysis}. Last rows:
PC variances $\hat\lambda_k=d_k^2/(n-1)$, $n=5$ days, and fractions of variance.}\label{ch07:tab-sixdays}
\begin{tabular}{@{}l rrrrr c rrr@{}}
\toprule
 & \multicolumn{5}{c}{$A$: centred daily changes (bp)} & & \multicolumn{3}{c}{loadings}\\
\cmidrule(lr){2-6}\cmidrule(l){8-10}
Tenor & Jan 03 & Jan 04 & Jan 05 & Jan 08 & Jan 09 & $s_j$ & PC1 & PC2 & PC3\\
\midrule
3M  & $-5.4$ & $-19.4$ & $-12.4$ & $19.6$ & $17.6$ & 17.70 & 0.38 & 0.53 & $-0.48$\\
6M  & $-4.6$ & $-14.6$ & $-12.6$ & $14.4$ & $17.4$ & 15.03 & 0.34 & 0.44 & $-0.05$\\
1Y  & $1.0$  & $-14.0$ & $-14.0$ & $9.0$  & $18.0$ & 14.12 & 0.36 & 0.26 & 0.23\\
2Y  & $9.6$  & $-10.4$ & $-16.4$ & $2.6$  & $14.6$ & 13.13 & 0.35 & $-0.03$ & 0.46\\
3Y  & $13.4$ & $-10.6$ & $-17.6$ & $1.4$  & $13.4$ & 13.99 & 0.37 & $-0.13$ & 0.43\\
5Y  & $18.6$ & $-11.4$ & $-15.4$ & $-0.4$ & $8.6$  & 14.03 & 0.35 & $-0.29$ & 0.12\\
7Y  & $20.8$ & $-11.2$ & $-14.2$ & $0.8$  & $3.8$  & 13.92 & 0.32 & $-0.37$ & $-0.23$\\
10Y & $20.8$ & $-12.2$ & $-11.2$ & $-0.2$ & $2.8$  & 13.37 & 0.30 & $-0.38$ & $-0.35$\\
20Y & $14.6$ & $-7.4$  & $-7.4$  & $0.6$  & $-0.4$ & 8.99  & 0.18 & $-0.28$ & $-0.36$\\
\midrule
\multicolumn{6}{@{}l}{$\hat\lambda_k$ (bp$^2$): 1323.9, 397.2, 34.2, 1.4, 0; total 1756.7} & & 75.4\% & 22.6\% & 1.9\%\\
\bottomrule
\end{tabular}
\end{table}

The singular values are $d=(72.77,\,39.86,\,11.70,\,2.39,\,0)$. Observations:
\begin{itemize}
  \item The total variance $1756.7$ equals the sum of the nine tenor variances $s_j^2$ (\cref{ch07:thm-pca}(iv)).
  \item PC5 has zero variance: five centred columns span at most four dimensions
        (\cref{ch07:prop-svd}(iv)). With real data one uses hundreds of days, yet the reading notes the
        ``curious fact'' that the loadings look almost the same as in the five-day toy.
  \item PC1 has positive, roughly equal loadings (a level shift); its standard deviation, 36.4 bp, exceeds
        that of every single tenor (at most 17.7 bp), because a unit-norm vector with nine positive entries
        of about $1/3$ adds up nine highly correlated moves.
  \item PC2 is short-minus-long (a change of slope), PC3 is negative at the two ends and positive in the
        middle (curvature; the overall sign is arbitrary). The daily scores (in bp) are $d_k\bm v_k$:
        \begin{align*}
          \text{PC1}&=(27.2,-37.6,-41.0,17.2,34.2)\ \text{(up, down, down, up, up)},\\
          \text{PC2}&=(-31.7,-4.5,3.1,18.5,14.6).
        \end{align*}
\end{itemize}

\begin{coursenote}[Erratum: dates in Table 1 of \file{pset_analysis}]
Table~1 of the reading labels the yield columns Jan 03, \dots, Jan 10 and the changes Jan 04, 05, 06, 07, 10;
but 6--7 January 2001 were a Saturday and a Sunday. The data are those of pages 5--6 of the reading: yields on
Jan 02, 03, 04, 05, 08, 09 and changes dated Jan 03, 04, 05, 08, 09, as used in \cref{ch07:tab-sixdays}.
\end{coursenote}

\subsection{Centring and scaling: covariance or correlation?}\label{ch07:sec-corr}
The typical workflow \ts{24}{9}{21:09}:
\begin{enumerate}[1.]
  \item transform the data to be as close as possible to multivariate
      normal (mean and covariance then describe the whole distribution: PCA only sees second moments);
  \item subtract the mean of each variable and possibly scale it to unit variance;
  \item compute the $p\times p$
      covariance (or correlation) matrix, symmetric and PSD by construction;
  \item diagonalise, order by decreasing
      eigenvalue, keep the first $K$ components.
\end{enumerate}

\paragraph{Centring is not optional.} PCA rotates about the origin \ts{24}{9}{25:28}. Feed it raw prices
(large positive numbers) and PC1 simply points from the origin to the centre of the cloud, capturing
nothing about co-movement \ts{24}{9}{26:31}. Slide 13 of \file{lec09} shows the same cloud before and after
centring. For daily yield changes the mean is tiny (a fraction of a basis point per day), so centring changes
little, but it remains good practice \ts{24}{9}{44:13}. In time series the \emph{rolling} mean can itself be
unstable, one more reason to work with changes rather than levels.

\paragraph{Scaling is a modelling choice} \ts{24}{9}{22:10}. Standardising $z_i=(x_i-\mu_i)/\sigma_i$ turns PCA of
$\Sigma$ into PCA of the correlation matrix $R=D_\sigma^{-1}\Sigma D_\sigma^{-1}$, $D_\sigma=\diag(\sigma_1,\dots,\sigma_p)$. Then
\begin{itemize}
  \item the total variance is $\tr R=p$ and every variable enters on the same footing;
  \item the eigenvectors of $R$ are \emph{not} the rescaled eigenvectors of $\Sigma$ (\cref{ch07:sec-geom}:
        PCA is not scale-equivariant);
  \item without scaling, a variable with much larger variance dominates PC1.
\end{itemize}
Andreev's toy example: the same data set looks one-dimensional (strongly elongated) in raw units and like
a structureless blob after normalising both axes \ts{24}{9}{23:12}. Rule of thumb \ts{24}{9}{24:24}: if the
coordinates have different units (temperature in Seattle versus pressure in Atlanta), standardise; if they
have the same units and their relative size is meaningful (grades on two quizzes, yield changes in bp),
covariance PCA keeps useful information about relative volatility. For Treasury yield changes Andreev uses
covariance PCA \ts{24}{9}{45:18}; how to normalise is ``very much the art'' and sometimes the research
question itself. For 2001--2005 Treasury changes the two choices give PC1/PC2/PC3 shares of
84.9/9.2/2.9\% (covariance) and 77.1/15.7/3.5\% (correlation); 2024 PS3 Problem~4 asks why.

For equity indices the slides use log returns $r_t=\ln(p_t/p_{t-1})$, which are closer to normal than price
changes; in that example covariance and correlation PCA give almost the same loadings because the index
variances are similar.

\subsection{How many components, and can we trust them?}\label{ch07:sec-robust}
PCA always returns an answer, even for pure noise \ts{24}{9}{12:44}. Andreev's advice: look at the
\emph{eigenvalues}, not only at the eigenvectors \ts{24}{9}{13:44}. A component is worth using if its eigenvalue
is well separated from the next one, ideally by an order of magnitude; with rolling PCA, plot the eigenvalue
(or variance-explained) time series on a log scale and check that the lines never cross \ts{24}{9}{14:50}. If
PC1 and PC2 swap places under small changes of the data, there is no stable structure and the PCA output
``obfuscates rather than clarifies'' \ts{24}{9}{15:54}. The standard display is the \emph{scree plot} of
$\hat\lambda_k$ against $k$ (\cref{ch07:fig-scree}).

\begin{intuition}[Degenerate perturbation theory (not discussed in class)]
Why do nearly equal eigenvalues make eigenvectors unstable? Estimation noise perturbs the covariance,
$\hat\Sigma=\Sigma+\delta\Sigma$ with $\delta\Sigma=O(N^{-1/2})$. First-order perturbation theory, exactly as for a
Hamiltonian in quantum mechanics, gives
\[
  \delta\lambda_k=\bm\gamma_k\T\,\delta\Sigma\,\bm\gamma_k,\qquad
  \delta\bm\gamma_k=\sum_{j\ne k}\frac{\bm\gamma_j\T\,\delta\Sigma\,\bm\gamma_k}{\lambda_k-\lambda_j}\,\bm\gamma_j .
\]
Eigenvalues are robust, but eigenvectors rotate by an amount inversely proportional to the \emph{gap}
$\lambda_k-\lambda_j$; for a degenerate level the eigenvectors are fixed by the perturbation, not by the data.
Andreev's ``move one point and watch PC1 swing'' demonstration \ts{24}{9}{11:38} is this formula at work.
\end{intuition}

\begin{figure}[ht]
\centering
\begin{tikzpicture}
\begin{axis}[width=0.72\linewidth,height=5.6cm,ymode=log,xlabel={component $k$},
  ylabel={fraction of variance},xmin=0.5,xmax=9.5,ymin=0.001,ymax=1.5,xtick={1,...,9},
  legend pos=north east,legend style={font=\small},grid=major,grid style={black!10},mark size=1.6pt]
  \addplot[thick,defcol,mark=*] coordinates {(1,0.8494) (2,0.0919) (3,0.0290) (4,0.0119) (5,0.00588)
    (6,0.00407) (7,0.00297) (8,0.00263) (9,0.00215)};
  \addlegendentry{covariance PCA}
  \addplot[thick,fincol,mark=square*] coordinates {(1,0.7708) (2,0.1571) (3,0.0355) (4,0.0175) (5,0.00806)
    (6,0.00481) (7,0.00245) (8,0.00207) (9,0.00167)};
  \addlegendentry{correlation PCA}
\end{axis}
\end{tikzpicture}
\caption{Scree plot (log scale) of the PCA of daily changes of 9 Treasury yields, 2001--2005 (2013 Case
Study~7; 2024 PS3 Problem~4). The first three eigenvalues are separated by roughly an order of magnitude
each, the pattern Andreev wants to see \ts{24}{9}{53:42}; standardising moves variance from PC1 to PC2.}\label{ch07:fig-scree}
\end{figure}

\paragraph{Effective dimension.} Slide 22 of \file{lec09} (after M.~Del Giudice, 2020) summarises a spectrum
in one number. With $\pi_k=\lambda_k/\sum_j\lambda_j$,
\[
  \mathrm{ED}=\frac{\bigl(\sum_k\lambda_k\bigr)^2}{\sum_k\lambda_k^2}=\frac{1}{\sum_k\pi_k^2}\in[1,p],
\]
the ``effective number of independent directions'' (the inverse participation ratio, familiar from
localisation physics): $\mathrm{ED}=1$ for a rank-one matrix such as the $4\times4$ matrix of ones, $2$ for two
decoupled blocks of perfectly correlated variables, $p$ for the identity. For the 2001--2005 Treasury covariance
PCA, ED $\approx1.37$.

\paragraph{Out-of-sample variance explained} \ts{24}{9}{20:07}. In finance PCA is calibrated on the past and
applied to the future. Let $W$ be the loadings estimated on a sample $X$ and let a new sample $Y$ have
covariance $\Sigma^Y=W^Y\Lambda^Y(W^Y)\T$. The PCs defined by $W$ have, in the new sample, covariance matrix
\begin{equation}\label{ch07:eq-oos}\begin{gathered}
  W\T\Sigma^YW=D\T\Lambda^YD,\qquad D=(W^Y)\T W\ \text{(orthogonal)},\\
  \text{share of PC }k=\frac{(D\T\Lambda^YD)_{kk}}{\tr\Lambda^Y}=\frac{\sum_jD_{jk}^2\lambda^Y_j}{\sum_j\lambda^Y_j}.\end{gathered}
\end{equation}
Since $\sum_jD_{jk}^2=1$, the out-of-sample variance of calibrated PC~$k$ is a weighted average of the new
eigenvalues: if the directions have rotated ($D\neq I$), variance leaks from the top factors to lower ones, and
by \cref{ch07:prop-best} the first $K$ calibrated PCs can never explain more than the first $K$ PCs of the new
sample itself. The off-diagonal entries of $D\T\Lambda^YD$ show that calibrated PCs are no longer uncorrelated
out of sample \ts{24}{9}{53:42}: in the Treasury example the correlations between next-day PC1, PC2, PC3 moves
are a few percent (slide 19: $7.1\%$ between PC1 and PC2), small enough for the factors to be called stable
\ts{24}{9}{54:47}.

\begin{coursenote}[A verbal slip]
At \ts{24}{9}{19:07} the speaker says that the ``product of eigenvalues'' is the variance of the data set; the
slide (and \cref{ch07:thm-pca}(iv)) has the correct statement: total variance $=\tr\Sigma=$ \emph{sum} of the
eigenvalues. (The product is $\det\Sigma$, the ``generalised variance''.)
\end{coursenote}

\begin{exercise}[not from the course]\label{ch07:ex-oos}
In \eqref{ch07:eq-oos} show that $\sum_k(D\T\Lambda^YD)_{kk}=\tr\Lambda^Y$ and that, for every $K$,
$\sum_{k\le K}(W\T\Sigma^YW)_{kk}\le\sum_{k\le K}\lambda^Y_k$. When does equality hold for $K=1$?
\end{exercise}

\subsection{PCA on financial time series}\label{ch07:sec-ts}
What is special about finance is that the data are \emph{time series} \ts{24}{9}{27:32}: each day is one
observation and the dimension is the number of variables observed that day. Consequences \ts{24}{9}{28:32}:
\begin{itemize}
  \item \textbf{Model changes, not levels} \ts{24}{9}{29:32}: levels (prices, yields) are non-stationary and not
        centred, while P\&L depends on changes after the bet is placed.
  \item \textbf{Hyperparameters} \ts{24}{9}{31:40}: the length of the look-back window, the weighting of old
        versus recent observations (rolling window, exponential weights), the set of variables included.
        More history or more variables is not always better: the marginal data may be noisier or less
        representative of where the bets will be placed.
  \item \textbf{Sampling frequency} \ts{24}{9}{33:46}: daily data are the default (long, clean histories);
        intraday data add observations but also microstructure noise. Volatility scales like
        $\sqrt{\text{time}}$, so the sampling interval sets its units: annualise daily numbers with $\sqrt{252}$
        before comparing \ts{24}{9}{34:51} (\cref{ch:volatility}).
  \item \textbf{Rolling PCA}: re-estimating every day turns loadings and eigenvalues into time series.
        Slowly varying loadings are reassuring; long stable periods interrupted by jumps signal regime
        changes \ts{24}{9}{30:33}.
\end{itemize}

% ------------------------------------------------------------------
\section{The Treasury yield curve: level, slope and curvature}\label{ch07:sec-yield}

\paragraph{The market.} The US Treasury funds the debt by regular auctions of bills and of 2, 3, 5, 7, 10,
(since 2020) 20 and 30-year notes and bonds \ts{24}{9}{38:00}; the most liquid points also trade as CME
futures (TU $\approx1.5$y, FV $\approx4$y, TY $\approx7$y, UXY $\approx10$y, US $\approx15$y, WN $\approx25$y equivalent maturities). The
market is the deepest in the world and allows large leverage \ts{24}{9}{35:55}; the Federal Reserve publishes
free, high-quality fitted curves \ts{24}{9}{36:58}. All these securities are loans to the same borrower, and
their prices depend on the expected path of the short rate set by the Fed: yield changes of different
maturities are therefore very highly correlated \ts{24}{9}{39:04}. In the 2024 data (par yields, daily changes)
the correlation is $98\%$ between 2y and 3y, $81\%$ between 2y and 10y, $68\%$ between 2y and 30y and $92\%$
between 10y and 30y (slide 17) \ts{24}{9}{40:07}: ``whenever you see a correlation matrix with high
concentrations of high correlations, it is a prime candidate for PCA''. The 2013 case study finds the
same pattern for 2001--2005: correlations above $0.9$ near the diagonal, decaying to $0.24$ between 3 months
and 20 years \ts{13}{15}{1:15:26}. Typical daily volatilities are 4--7 bp \ts{13}{15}{1:14:24}.

\paragraph{The yield curve} (\cref{sec:bonds}, \cref{fig:yieldcurve}) is usually upward sloping but can be
inverted, as in 2023--2024 \ts{24}{9}{41:11}; the Fed sets its short end (the 50 bp cut of September 2024
moved the leftmost point), the rest is market expectation \ts{24}{9}{43:12}. A movie of the curve would show
``strings flopping around'': mostly moving up and down, a little tilting, a little bending \ts{24}{9}{42:12}.

\paragraph{The three factors.} Running PCA on daily yield changes gives, in every period and data set
(\cref{ch07:fig-loadings}, \cref{ch07:tab-share}):
\begin{itemize}
  \item \textbf{PC1 = level}: all loadings of the same sign; yields move up or down together. It explains
        roughly 85--95\% of the variance \ts{24}{9}{47:24}, \ts{13}{15}{1:19:47}.
  \item \textbf{PC2 = slope}: loadings negative at the short end and positive at the long end (or vice versa):
        steepening versus flattening; 5--10\% \ts{24}{9}{48:24}.
  \item \textbf{PC3 = curvature}: same sign at both ends, opposite sign in the ``belly'': a second difference;
        1--3\% \ts{13}{15}{1:20:51}.
\end{itemize}
The top three factors explain 97--99\% of the daily variance.

\begin{figure}[ht]
\centering
\begin{tikzpicture}
\begin{axis}[width=0.52\linewidth,height=5.4cm,xmode=log,xmin=0.2,xmax=25,ymin=-0.6,ymax=0.6,
  xtick={0.25,0.5,1,2,5,10,20},xticklabels={3m,6m,1y,2y,5y,10y,20y},
  xlabel={tenor},ylabel={loading},title={\small (a) CMT yields, 2001--2005},
  grid=major,grid style={black!10},mark size=1.3pt,tick label style={font=\scriptsize},
  legend style={font=\scriptsize,at={(0.03,0.03)},anchor=south west}]
  \addplot[thick,defcol,mark=*] coordinates {(0.25,0.1088) (0.5,0.1581) (1,0.2531) (2,0.3905) (3,0.4197)
    (5,0.4206) (7,0.4009) (10,0.3697) (20,0.3101)};
  \addlegendentry{PC1 level}
  \addplot[thick,thmcol,mark=square*] coordinates {(0.25,-0.5225) (0.5,-0.4817) (1,-0.4107) (2,-0.2071)
    (3,-0.0873) (5,0.1082) (7,0.2289) (10,0.2837) (20,0.3621)};
  \addlegendentry{PC2 slope}
  \addplot[thick,intcol,mark=triangle*] coordinates {(0.25,0.5434) (0.5,0.2508) (1,-0.0514) (2,-0.4750)
    (3,-0.4022) (5,-0.0311) (7,0.1475) (10,0.2710) (20,0.3944)};
  \addlegendentry{PC3 curvature}
  \addplot[black!40,forget plot] coordinates {(0.2,0) (25,0)};
\end{axis}
\end{tikzpicture}\hfill
\begin{tikzpicture}
\begin{axis}[width=0.52\linewidth,height=5.4cm,xmin=0,xmax=31,ymin=-0.45,ymax=0.45,
  xlabel={maturity (years)},title={\small (b) par yields, Sep 2021--Aug 2023},
  grid=major,grid style={black!10},mark size=0.9pt,tick label style={font=\scriptsize}]
  \addplot[thick,defcol,mark=*] coordinates {(1,0.127) (2,0.185) (3,0.202) (4,0.208) (5,0.211) (6,0.212)
    (7,0.211) (8,0.209) (9,0.206) (10,0.203) (11,0.199) (12,0.195) (13,0.192) (14,0.188) (15,0.185) (16,0.181)
    (17,0.178) (18,0.176) (19,0.173) (20,0.171) (21,0.169) (22,0.167) (23,0.166) (24,0.164) (25,0.163)
    (26,0.162) (27,0.161) (28,0.160) (29,0.160) (30,0.159)};
  \addplot[thick,thmcol,mark=square*] coordinates {(1,-0.357) (2,-0.420) (3,-0.373) (4,-0.307) (5,-0.241)
    (6,-0.181) (7,-0.130) (8,-0.086) (9,-0.049) (10,-0.017) (11,0.010) (12,0.033) (13,0.052) (14,0.069)
    (15,0.084) (16,0.096) (17,0.107) (18,0.117) (19,0.126) (20,0.133) (21,0.140) (22,0.146) (23,0.151)
    (24,0.156) (25,0.161) (26,0.165) (27,0.168) (28,0.172) (29,0.175) (30,0.177)};
  \addplot[thick,intcol,mark=triangle*] coordinates {(1,0.273) (2,0.286) (3,0.260) (4,0.115) (5,-0.050)
    (6,-0.185) (7,-0.274) (8,-0.317) (9,-0.321) (10,-0.298) (11,-0.257) (12,-0.204) (13,-0.146) (14,-0.088)
    (15,-0.033) (16,0.017) (17,0.061) (18,0.097) (19,0.125) (20,0.146) (21,0.160) (22,0.166) (23,0.166)
    (24,0.160) (25,0.149) (26,0.133) (27,0.112) (28,0.087) (29,0.059) (30,0.028)};
  \addplot[black!40] coordinates {(0,0) (31,0)};
\end{axis}
\end{tikzpicture}
\caption{Loadings of the first three PCs of daily yield changes (covariance PCA). (a) 9 FRED constant-maturity
yields, 2001--2005, \texttt{princomp} output of 2013 Case Study~7. (b) Fed par yields at maturities 1--30 years,
504-day window ending 1 September 2023, computed from \file{fedbondyielddata_20230911.csv} as in
\file{feddataanalysis.ipynb} (signs fixed so that the 20--30y loadings sum to a positive number); PC1--PC3
explain 89.2\%, 8.3\% and 1.0\%. Compare slide 18 of \file{lec09}.}\label{ch07:fig-loadings}
\end{figure}

\begin{table}[ht]
\centering\small
\caption{Fraction of the variance of daily yield changes explained by the first three PCs.}\label{ch07:tab-share}
\setlength{\tabcolsep}{3pt}\begin{tabular}{@{}l l rrrr@{}}
\toprule
Data & Source & PC1 & PC2 & PC3 & top 3\\
\midrule
9 CMT yields, 6 days of 2001 & 2016 reading & 75.4 & 22.6 & 1.9 & 99.9\\
9 CMT yields, 2001--2005, covariance & Case Study 7 & 84.9 & 9.2 & 2.9 & 97.0\\
9 CMT yields, 2001--2005, correlation & PS7 (2013), PS3 (2024) & 77.1 & 15.7 & 3.5 & 96.3\\
9 CMT yields, 2009--2013, covariance & Case Study 7 & 88.6 & 6.8 & 1.6 & 96.9\\
30 par yields, 2-year window to Sep 2023 & computed from the 2024 data & 89.2 & 8.3 & 1.0 & 98.5\\
\bottomrule
\end{tabular}
\end{table}

\begin{intuition}[Modes of a vibrating string]
Level, slope and curvature are the analogues of the first normal modes of a string: 0, 1 and 2 sign changes
(nodes). This is partly a mathematical consequence of any smooth, strongly and positively correlated
curve (not discussed in class): for the correlation kernel $\rho_{ij}=e^{-\kappa|\tau_i-\tau_j|}$ on the same nine
tenors with $\kappa=0.1$, the first four eigenvectors have exactly 0, 1, 2, 3 sign changes and explain
66\%, 15\%, 9\%, 4\% (my computation; such kernels are ``oscillation matrices'' in the sense of
Gantmacher and Krein). The economics lies in the \emph{details}: the precise shapes, the explained
variances and how they change over time. The same holds for any Karhunen--Lo\`eve expansion of a random
field with a smooth covariance.
\end{intuition}

\paragraph{Reading the loadings as regime indicators} \ts{24}{9}{49:29}. The PC1 loading is small at the very
short end and levels off beyond a few years: short rates are anchored by the Fed, long rates by
expectations. The maturity at which PC1 ``levels off'' (in the notebook: the first tenor with loading above
0.15) was 1--3 years until the financial crisis, rose to 5--7 years when the Fed promised to keep rates
low for long (quantitative easing, forward guidance), fell back, jumped again with the COVID response and
decreased during the 2022 tightening \ts{24}{9}{50:33}, \ts{24}{9}{1:14:43}. Similarly the 2-year PC1 loading
collapsed after 2012 (the Fed was not moving) and rebounded when policy became uncertain; the
\emph{location of the belly} of PC3 (usually near 10y in the US) shifts with regimes \ts{24}{9}{51:34}.
Andreev's macro reading of curvature \ts{24}{9}{1:17:51}: the front end is driven by the central bank, the
intermediate part by inflation expectations, the long end by growth expectations; investors also differ by
segment (speculators and short-term funding at 2y, corporate issuance at 5--10y, pension funds and insurers
matching long liabilities at the long end) \ts{24}{9}{1:18:57}. In an emerging market such as Mexico the belly
sits nearer 5 years, reflecting shorter investment horizons. Once the level and slope are hedged out, these
small but structured features emerge from ``the big noise of rates going up and down'' \ts{24}{9}{1:19:59}.

\paragraph{The last PC is interesting too} \ts{13}{15}{1:21:57}. In the 2013 case study PC9, with the smallest
variance (0.8 bp/day), is essentially the 7-year yield against a weighted average of the 5- and 10-year: the
most stable linear combination of yields, i.e.\ the best hedge of the 7-year point with its neighbours.
Cumulating the scores shows the history of the factors \ts{13}{15}{1:23:00}: over 2001--2005 cumulative PC1
fell from 0 to about $-8$ and came back (rates fell, then rose), cumulative PC2 rose to about 6 and fell back
(the curve steepened, then flattened), curvature moved much less \ts{13}{15}{1:24:00}.

\begin{finance}[Nelson--Siegel--Svensson curves (not discussed in class)]
The Fed file \file{fedbondyielddata_20230911.csv} is the Gurkaynak--Sack--Wright data set: daily zero-coupon
yields (\texttt{SVENY01}--\texttt{30}, continuously compounded), par yields (\texttt{SVENPY}), instantaneous and
one-year forward rates (\texttt{SVENF}, \texttt{SVEN1F}) from 1961, and the parameters \texttt{BETA0}--\texttt{BETA3},
\texttt{TAU1}, \texttt{TAU2} of the fitted Svensson curve
\[
  y(n)=\beta_0+\beta_1\frac{1-e^{-n/\tau_1}}{n/\tau_1}
        +\beta_2\Bigl[\frac{1-e^{-n/\tau_1}}{n/\tau_1}-e^{-n/\tau_1}\Bigr]
        +\beta_3\Bigl[\frac{1-e^{-n/\tau_2}}{n/\tau_2}-e^{-n/\tau_2}\Bigr]
\]
(I checked that this reproduces \texttt{SVENY} to four decimals). The three Nelson--Siegel basis functions
are a parametric version of level ($\beta_0$), slope ($\beta_1$, which dies out with maturity) and curvature
($\beta_2$, a hump); $\beta_3$ adds a second hump. Because the published curves are smooth fits, their PCA
loadings are smoother than those of raw bond yields.
\end{finance}

\begin{casestudy}[Reading \file{pset_analysis} and its R code \file{pset_rcode}: five years of yields]
\textbf{Goal.} Reproduce the six-day example and extend it to 2001--2005. \textbf{Data.} FRED
constant-maturity yields \texttt{DGS3MO}, \dots, \texttt{DGS20} (script \file{fm_freddata.r}), 1247 daily changes.
\textbf{Method.} \texttt{svd()} of the centred days$\times$tenors matrix, $\hat\lambda_k=d_k^2/(n-1)$, bar plots of
loadings, daily and cumulative PC1/PC2 series, \texttt{pairs()} and interactive 3D plots of (PC1, PC2, PC3);
then the built-in \texttt{princomp()} and \texttt{screeplot()}, which give the same result up to the signs of
the loadings (\cref{ch07:rem-sign}). \textbf{Results.} Total variance 0.0306 (\%$^2$); PC1/PC2/PC3 explain
84.8/9.3/2.9\%; PC1 loadings 0.11 (3m) rising to 0.42 (3y--5y) and 0.31 (20y); cumulative PC1 falls
over the first half of the period and recovers, cumulative PC2 (the spread) rises then falls.
\textbf{Lesson.} Five days already reveal the structure that five years confirm.
\emph{Two small inconsistencies:} in the five-year section the text says $n$ is ``the number of columns
(days)'' but the days are now rows ($\hat\lambda_k=d_k^2/(\texttt{NROW(X)}-1)$ in the code); and the text calls PC2
``shorter minus longer'' although in this SVD its loadings are negative at the short end, i.e.\ long minus
short. Signs are arbitrary.
\end{casestudy}

\begin{casestudy}[\file{feddataanalysis.ipynb}: Andreev's notebook on the Fed data]
\textbf{Data.} Par yields \texttt{SVENPY01}--\texttt{30} from the first day with a 30-year yield (25 Nov 1985) to
1 Sep 2023, 9430 days, daily differences. \textbf{Method} (Python/pandas): 504-day (2-year) rolling covariance of
yield changes, \texttt{numpy.linalg.eigh} each day, signs normalised so that the 20--30y loadings sum to a
positive number; variance explained per factor (area and log plots); PC1 loading stability for 2, 5, 10, 30y;
``maturity for PC1 leveling'' (first tenor with loading $>0.15$), PC2 pivot and PC3 belly; next-day yield changes
projected on today's loadings (the out-of-sample check \eqref{ch07:eq-oos}); out-of-sample variance explained with
loadings two years old; a PC1/PC2-neutral 5s7s10s fly and a mean-reversion strategy on it.
\textbf{Results.} Correlation matrix of slide 17 (reproduced exactly); out-of-sample correlations between
next-day PC1, PC2, PC3 moves: $4.0\%$ (PC1/PC2), $-0.2\%$ (PC1/PC3), $-9.0\%$ (PC2/PC3) (the slide shows
$7.1\%$, $-1.6\%$, $-7.3\%$ from another run); strategy Sharpe ratio $0.10$ (slide: $0.09$).
\textbf{Lesson.} The last cell repeats the strategy with the signal \emph{not} delayed by one day and gets
spectacular, unrealistic returns: a textbook demonstration of look-ahead bias. Andreev recommends the
notebook as a template for a final project \ts{24}{9}{1:11:29}.
\end{casestudy}

% ------------------------------------------------------------------
\section{PC portfolios: curves, flies and hedging}\label{ch07:sec-pcport}

\subsection{Loadings as portfolio weights}
Why care about PC2 and PC3 if PC1 explains 90\% of the variance? Because they are where relative-value
trading happens \ts{24}{9}{55:49}. The six liquid tenors are so correlated that trading any of them outright
is essentially a bet on PC1 \ts{24}{9}{56:51}; PCA manufactures from them a handful of \emph{uncorrelated}
portfolios, i.e.\ more independent bets \ts{24}{9}{57:54}.

Positions in bonds are measured in \emph{dv01}: \$1 dv01 of a bond is the notional whose value changes by
\$1 for a 1 bp change of its yield (\cref{sec:duration}). Let $h_i$ be the dv01 held in tenor $i$ (positive =
long the bond, which gains when the yield falls). To first order the daily P\&L in dollars is
\[
  \mathrm{PnL}=-\bm h\T\Delta\bm y=-\sum_k(\bm h\T\bm\gamma_k)\,p_k,\qquad \Delta\bm y\ \text{in bp},
\]
using $\Delta\bm y=\sum_kp_k\bm\gamma_k$ (\cref{ch07:thm-pca}(iii), centring neglected). The \emph{exposure} of the portfolio to
PC~$k$ is $e_k=\bm h\T\bm\gamma_k$. Buying ``the PC1 portfolio'' means holding every bond in proportion to its PC1
loading \ts{24}{9}{58:59}; a portfolio with $e_1=0$ is \emph{PC1-neutral}, one with $e_1=e_2=0$ is PC1/PC2-neutral.

\begin{coursenote}[Erratum: the dv01 box on slide 20 of \file{lec09}]
The box defines ``buy \$1 dv01 of bond'' as the notional such that the position \emph{earns} \$1 when the yield
goes \emph{up} 1 bp. A long bond position loses when yields rise; the examples on the same slide (a steepener
sells the 10-year bond, a belly fly buys the 10-year) use the standard convention adopted here. The notebook
avoids the issue by computing returns directly in yield space ($+$weights$\cdot\Delta\bm y$), i.e.\ as positions in
yields rather than in bonds; any sign convention works if used consistently.
\end{coursenote}

\paragraph{PC-neutral curves} (``PC1 residuals''). Given two tenors $a,b$, the portfolio $\bm h=\bm e_a-\alpha\bm e_b$
(\$1 dv01 of $a$ against $\alpha$ dv01 of $b$) is PC1-neutral iff
\begin{equation}\label{ch07:eq-curve}
  \gamma_{a1}-\alpha\gamma_{b1}=0\iff\alpha=\frac{\gamma_{a1}}{\gamma_{b1}},
\end{equation}
where $\gamma_{ik}$ is the loading of tenor $i$ on PC~$k$. Example: the 2s10s \emph{steepener}, short \$1 dv01 of
10-year, long $\alpha$ dv01 of 2-year: it profits when the curve steepens and is insensitive to parallel (PC1) moves.
Such curve trades exist in the market as standard products \ts{24}{9}{1:00:00}.

\paragraph{PC-neutral flies} (``PC2 residuals''). Given three tenors, $\bm h=\alpha_a\bm e_a+\alpha_b\bm e_b-\bm e_c$ is neutral
to PC1 and PC2 iff
\begin{equation}\label{ch07:eq-fly}
  \begin{pmatrix}\gamma_{a1}&\gamma_{b1}\\ \gamma_{a2}&\gamma_{b2}\end{pmatrix}
  \begin{pmatrix}\alpha_a\\ \alpha_b\end{pmatrix}=\begin{pmatrix}\gamma_{c1}\\ \gamma_{c2}\end{pmatrix}.
\end{equation}
Example: the 2s10s30s ``belly'' fly buys \$1 dv01 of 10-year and sells $\alpha_2$ of 2-year and $\alpha_{30}$ of
30-year: a pure bet on curvature.

\begin{coursenote}[Erratum: slide 20 of \file{lec09}]
The fly system on the slide has $W_{b,0}$ in the lower-right entry; it must be $W_{b,1}$ (the PC2 loading of
tenor $b$), as in \eqref{ch07:eq-fly}. (The slides index PCs from 0.)
\end{coursenote}

\begin{example}[Hedge ratios and volatilities from the 2001--2005 PCA]\label{ch07:ex-hedges}
Rebuild $\Sigma=\Gamma\Lambda\Gamma\T$ from the loadings and standard deviations of 2013 Case Study~7 (this reproduces the
case-study volatilities and correlations to three digits) and measure risk as the standard deviation of daily P\&L
per \$1 dv01 of the reference bond:
\begin{center}\small
\begin{tabular}{@{}l l c c@{}}
\toprule
Trade (PC-neutrality) & Weights (dv01) & vol (\$/day) & outright vol\\
\midrule
2s10s curve (PC1) & $-1$ of 10y, $+0.947$ of 2y & 3.5 & 6.3 (10y)\\
2s5s10s fly (PC1, PC2) & $+1$ of 5y, $-0.423$ of 2y, $-0.690$ of 10y & 1.5 & 6.9 (5y)\\
5s7s10s fly (PC1, PC2) & $+1$ of 7y, $-0.367$ of 5y, $-0.667$ of 10y & 1.0 & 6.7 (7y)\\
\bottomrule
\end{tabular}
\end{center}
The 5s7s10s fly is the notebook's trade and is close to PC9 of the case study. To run it at the risk of an
outright 7-year position one needs about $6.4$ times the dv01 on the belly, plus the offsetting wings.
\end{example}

\subsection{Leverage, and why higher PCs are the interesting ones}
A PC2 portfolio moves perhaps a third as much as the PC1 portfolio per unit notional; scaling it up three
times gives the same daily standard deviation and a return stream uncorrelated with the market direction
\ts{24}{9}{1:01:03}. Far from being irrelevant, the higher PCs are ``where a lot of the money is made''
\ts{24}{9}{1:03:07}; beyond PC3, however, the leverage needed becomes too expensive. The strategy requires a
liquid market and cheap financing (repo); regulators (the CFTC and others) watch the resulting leveraged
long/short Treasury positions as a possible systemic risk \ts{24}{9}{1:02:05}. In equities PC1 is much less
dominant, so less leverage is needed; market-neutral long/short equity funds are precisely ``hedging out PC1''
\ts{24}{9}{1:04:09}.

\begin{center}\small
\begin{tabularx}{\linewidth}{@{}X X@{}}
\toprule
\textbf{Pros of PC-residual portfolios} & \textbf{Cons} (slide 21)\\
\midrule
Bets unrelated to the main market move (PC1). & Transaction cost per unit of volatility is much higher.\\
More effectively uncorrelated bets (diversity). & Weights are model-dependent and change over time.\\
Low volatility per notional: a mispricing may have a better signal-to-noise ratio. &
Much larger notionals, i.e.\ leverage and balance sheet; financing (overnight repo) can dry up in stress.\\
\bottomrule
\end{tabularx}
\end{center}

\paragraph{Model risk.} The hedge ratios are only as good as the loadings. When COVID hit and the Fed made an
emergency cut, the curve dynamics changed overnight \ts{24}{9}{1:07:21}; a leveraged PC2 position computed with
stale loadings is no longer PC1-neutral and its P\&L is dominated by unintended level risk, while there is no
history yet to recalibrate on \ts{24}{9}{1:08:21}. Possible answers: implied (option) markets, or calibrating
on past periods of a similar regime (e.g.\ all Fed cutting cycles), with its own pitfalls \ts{24}{9}{1:09:25}.
Event days also matter for risk: daily yield volatility is typically 5--6 bp, but the market may price twice
that on the day of a major inflation release \ts{24}{9}{1:06:18}.

\paragraph{Diversity of PC portfolios.} Slide 23 computes the effective dimension of the correlation matrix of
the returns of PC portfolios: PC1 alone 1, PC1+PC2 1.99, PC1+PC2+PC3 2.97 (almost exactly uncorrelated also
out of sample), and 5.28 when the residual portfolios left after removing three PCs are added. Higher-order
residuals also tend to \emph{mean-revert}. The notebook's strategy \ts{24}{9}{1:10:25}: build the PC1/PC2-neutral
5s7s10s fly every day from the current loadings; signal $s$ = z-score of its cumulative return against an
exponentially weighted mean (half-life 100 days) and standard deviation (half-life 60 days), delayed by one day;
position $=-s$. The Sharpe ratio is only about 0.1: the point is the method, not the strategy.

\begin{coursenote}[Differences between slides and speech]
At \ts{24}{9}{1:10:25} the speaker describes the strategy as ``momentum/mean reversion for PC1''; slide 24 and the
notebook apply mean reversion to the PC1/PC2-neutral fly.
\end{coursenote}

\subsection{Hedging a portfolio against PC1}
For an equity portfolio with weights $\bm V$ on $p$ indices (returns, so $\mathrm{PnL}=\bm V\T\bm r$), the exposure to
PC~$n$ is $\bm V\T\bm\gamma_n$. Projecting out the first component,
\[
  \bm V'=\bm V-(\bm V\T\bm\gamma_1)\,\bm\gamma_1,\qquad \bm V'^{\mathsf T}\bm\gamma_1=0,\qquad
  \Var(\bm V'^{\mathsf T}\bm r)=\bm V'^{\mathsf T}\Sigma\bm V'=\bm V\T\Sigma\bm V-\lambda_1(\bm V\T\bm\gamma_1)^2,
\]
the last identity because $(I-\bm\gamma_1\bm\gamma_1\T)\Sigma(I-\bm\gamma_1\bm\gamma_1\T)=\Sigma-\lambda_1\bm\gamma_1\bm\gamma_1\T$. If the expected return
(``alpha'') of the strategy is roughly orthogonal to PC1, hedging keeps the alpha and removes most of the
risk: higher Sharpe ratio.

\begin{coursenote}[Erratum: slide 32 of \file{lec09}]
``The stdev of the resulting portfolio $=V'^T\Sigma V'$'': that is the variance; the standard deviation is
$\sqrt{V'^T\Sigma V'}$.
\end{coursenote}

\begin{casestudy}[European stock indices (slides 26--34; not covered in class \ts{24}{9}{1:12:30})]
\textbf{Data.} Daily closes of AEX, DAX, FTSE MIB, IBEX, OMX, Euro Stoxx 50 and CAC, 4 Jan 2010 to 16 Sep 2016
(same time zone and currency), rebased to 100 EUR and converted to demeaned log returns. Correlations 72--98\%.
\textbf{PCA.} PC1 explains 89\% with nearly equal loadings: a ``Euro stocks'' factor. PC2 (about 5\%, half of the
rest) separates North (Netherlands, Germany, Sweden: positive) from South (Italy, Spain: negative), France near
zero. Covariance and correlation PCA agree. Rolling 6-month PCA: top three factors explain 90--98\% throughout,
PC1 loadings stable, PC2 fairly stable. \textbf{Hedging example} (thesis: German stocks outperform, 2010, target
100m EUR annualised volatility): buying the DAX (vol 20\%) needs 500m and returned 75m (Sharpe 0.75); the
PC1-neutral portfolio $\bm e_{\mathrm{DAX}}-0.377\,\bm\gamma_1$ (weights $+0.86$ DAX, about $-0.14$ each other index) has vol
6.7\%, needs 1{,}500m and returned about 230m (Sharpe 2.3). \textbf{Other uses:} North-versus-South trades;
Frankfurt opens before Madrid, so a PCA model can price Spain by proxy just before its open.
\emph{Slide typos:} the rebasing date is 4 Jan 2010 (not ``1/4/2016''), and the legend entry ``Swiss'' on slides 27
and 31 is the Euro Stoxx 50.
\end{casestudy}

\begin{finance}[Limitations (slide 35)]
PCA always gives an answer, so test robustness to the window and to the set of variables; it works best on data
close to multivariate normal; and it implies no causality. To predict a dependent variable, use regression
(\cref{ch:regression}). Andreev's closing advice \ts{24}{9}{1:22:03}: data that are hard to get (illiquid markets,
emerging-market curves) are where fewer competitors look, hence where quantitative modelling may still pay.
\end{finance}

% ------------------------------------------------------------------
\section{Linear factor models}\label{ch07:sec-lfm}

\subsection{Definition and covariance structure}
\begin{definition}[Linear factor model]\label{ch07:def-lfm}
For $m$ assets observed over $T$ periods \ts{13}{15}{1:06}, \ts{13}{15}{3:13}:
\begin{equation}\label{ch07:eq-lfm}
  x_{i,t}=\alpha_i+\beta_{1,i}f_{1,t}+\dots+\beta_{K,i}f_{K,t}+\epsilon_{i,t},\qquad
  \bm x_t=\bm\alpha+B\bm f_t+\bm\epsilon_t ,
\end{equation}
with intercepts $\bm\alpha\in\R^m$, \emph{factor loadings} $B\in\R^{m\times K}$ (row $i$ is $\bm\beta_i\T$), \emph{common factors}
$\bm f_t\in\R^K$, $K\ll m$, and \emph{specific factors} $\bm\epsilon_t$. Assumptions \ts{13}{15}{5:19}:
$\{\bm f_t\}$ is covariance stationary with $\E[\bm f_t]=\bm\mu_f$, $\Cov\bm f_t=\Omega_f$; $\{\bm\epsilon_t\}$ is white noise with
$\E[\bm\epsilon_t]=\bm0$, $\Cov\bm\epsilon_t=\Psi=\diag(\sigma_1^2,\dots,\sigma_m^2)$, $\Cov(\bm\epsilon_t,\bm\epsilon_{t'})=0$ for $t\neq t'$; and
$\Cov(\bm f_t,\bm\epsilon_{t'})=0$ for all $t,t'$.
\end{definition}

The key assumption is that $\Psi$ is \emph{diagonal}: all co-movement between assets goes through the factors.

\begin{proposition}[Moments of the factor model]\label{ch07:prop-lfm}
$\E[\bm x_t\mid\bm f_t]=\bm\alpha+B\bm f_t$, $\Cov[\bm x_t\mid\bm f_t]=\Psi$, and unconditionally
\begin{equation}\label{ch07:eq-lfmcov}
  \E[\bm x_t]=\bm\alpha+B\bm\mu_f,\qquad \Sigma_x=\Cov[\bm x_t]=B\,\Omega_f\,B\T+\Psi .
\end{equation}
\end{proposition}
\begin{proof}
By the law of total covariance, $\Cov[\bm x_t]=\E\bigl[\Cov(\bm x_t\mid\bm f_t)\bigr]+\Cov\bigl(\E[\bm x_t\mid\bm f_t]\bigr)=\Psi+B\Omega_fB\T$;
directly, $\Cov(B\bm f_t+\bm\epsilon_t)=B\Omega_fB\T+\Psi+2\Cov(B\bm f_t,\bm\epsilon_t)$ and the cross term vanishes \ts{13}{15}{7:25}.
\end{proof}

\paragraph{Parameter count} \ts{13}{15}{8:27}. An unrestricted covariance has $m(m+1)/2$ parameters; the factor
structure has $mK$ (loadings) $+K(K+1)/2$ ($\Omega_f$) $+m$ ($\Psi$) \ts{13}{15}{10:39}. For $m=500$ stocks and $K=10$
factors: $125{,}250$ versus $5{,}555$. Moreover each new asset adds only $K+1$ parameters. This is why factor models
are used for portfolio optimisation, where the (inverse) covariance matrix drives the allocation
(\cref{ch:portfolio}), and for risk management, where scenarios are naturally expressed as factor moves
\ts{13}{15}{16:54}.

\paragraph{Three regression views} (\cref{ch:regression}).
\emph{Cross-sectional}: fix $t$; \eqref{ch07:eq-lfm} is a regression of the $m$ returns on the columns of $B$ with
``coefficients'' $\bm f_t$ \ts{13}{15}{4:18}. \emph{Time series}: fix asset $i$; $\bm x_i=\bm1_T\alpha_i+F\bm\beta_i+\bm\epsilon_i$
with $F\in\R^{T\times K}$ the factor history and $\Cov\bm\epsilon_i=\sigma_i^2I_T$: an ordinary regression satisfying the
Gauss--Markov assumptions \ts{13}{15}{11:41}. \emph{Multivariate}: stacking all assets,
$X=\bm1_T\bm\alpha\T+FB\T+E$ with $X=[\bm x_1|\cdots|\bm x_m]\in\R^{T\times m}$ \ts{13}{15}{12:45} (the slides write $B$ for the
$K\times m$ matrix of the time-series form, the transpose of the cross-sectional $B$).
Whether factors or loadings are observed defines the three families below \ts{13}{15}{0:00}.

\subsection{Identification}\label{ch07:sec-ident}
When the factors are latent, the model is not identified \ts{13}{15}{42:40}: for any invertible $H\in\R^{K\times K}$,
$\bm f_t^*=H\bm f_t$ and $B^*=BH^{-1}$ give the same $\bm x_t$, with $\bm\mu_f^*=H\bm\mu_f$ and $\Omega_f^*=H\Omega_fH\T$
\ts{13}{15}{43:49}. One uses this freedom to normalise \ts{13}{15}{44:50}:
\begin{itemize}
  \item \emph{orthonormal factors}: if $\Omega_f=\Gamma_f\Lambda_f\Gamma_f\T$, take $H=\Lambda_f^{-1/2}\Gamma_f\T$, so that
        \[\Omega_f^*=\Lambda_f^{-1/2}\Gamma_f\T\Gamma_f\Lambda_f\Gamma_f\T\Gamma_f\Lambda_f^{-1/2}=I_K;\]
  \item \emph{zero-mean factors}: absorb $B\bm\mu_f$ into the intercept, $\bm\alpha^*=\bm\alpha+B\bm\mu_f$.
\end{itemize}
Then $\Sigma_x=BB\T+\Psi$ \ts{13}{15}{46:00}. An indeterminacy remains: $B\mapsto BQ$ with $Q$ orthogonal leaves $BB\T$
unchanged. This rotational freedom is exploited to make factors interpretable (\cref{ch07:sec-fa}).

\begin{coursenote}[Erratum: slide 22 of \file{lecnote15}]
The slide takes $H=\Gamma\Lambda^{-1/2}$; then $H\Omega_fH\T=\Gamma\Lambda^{-1/2}\Gamma\Lambda\Gamma\T\Lambda^{-1/2}\Gamma\T$, which is not $I_K$ in general.
The correct choice is $H=\Lambda^{-1/2}\Gamma\T$ (or $Q\Lambda^{-1/2}\Gamma\T$ for any orthogonal $Q$).
\end{coursenote}

% ------------------------------------------------------------------
\section{Macroeconomic and fundamental factor models}\label{ch07:sec-macrofund}

\subsection{Observed factors: Sharpe's single-index model}
The simplest factor model \ts{13}{15}{13:46}: excess returns regressed on the excess return $R_{M,t}$ of the
market index,
\[
  x_{i,t}=\alpha_i+\beta_iR_{M,t}+\epsilon_{i,t},\qquad
  \Sigma_x=\sigma_M^2\,\bm\beta\bm\beta\T+\Psi,\quad \sigma_M^2=\Var(R_{M,t}) .
\]
The covariance between two different assets is $\beta_i\beta_j\sigma_M^2$ \ts{13}{15}{14:48}: a rank-one matrix plus a
diagonal, trivial to build even for thousands of securities \ts{13}{15}{15:53}. $\beta_i$ measures the systematic
risk of asset $i$ (the CAPM of \cref{ch:portfolio} adds the equilibrium restriction $\alpha_i=0$).

\paragraph{Estimation} \ts{13}{15}{17:57}. The time-series regression of each asset satisfies Gauss--Markov, so OLS
gives BLUE estimates $(\hat\alpha_i,\hat\beta_i)$ (and the MLE under normality); unbiased residual variances are
$\hat\sigma_i^2=\hat{\bm\epsilon}_i\T\hat{\bm\epsilon}_i/(T-2)$; $\hat\sigma_M^2$ is the sample variance of $R_M$; and
$\hat\Sigma_x=\hat\sigma_M^2\hat{\bm\beta}\hat{\bm\beta}\T+\hat\Psi$.

\begin{intuition}[What PCA sees in a single-index world (not discussed in class)]
If all specific variances are equal, $\Psi=\psi I$, then $\Sigma_x=\sigma_M^2\bm\beta\bm\beta\T+\psi I$ has top eigenvector
$\bm\beta/\|\bm\beta\|$ with eigenvalue $\sigma_M^2\|\bm\beta\|^2+\psi$, and all other eigenvalues equal to $\psi$. PCA recovers the
market factor exactly and finds nothing else (a degenerate remainder, \cref{ch07:rem-sign}). With unequal
$\sigma_i^2$ the recovery is only approximate: this is the difference between PCA and factor analysis
(\cref{ch07:sec-pcafa}).
\end{intuition}

\subsection{Macroeconomic multifactor models}
Any observable series is a candidate factor, judged by how much it contributes to explaining returns and
covariances \ts{13}{15}{19:00}: market risk, inflation (CPI, PPI, commodity prices), industrial production, money
growth, interest rates, housing starts, unemployment. Chen, Roll and Ross (1986) used not the levels but the
\emph{surprises}: unanticipated changes, computed as innovations of a vector autoregression fitted to the macro
variables (\cref{ch:timeseries}) \ts{13}{15}{20:03}, \ts{13}{15}{21:07}. Since $F$ is observed, each asset's
loadings are again time-series OLS estimates, with $\hat\sigma_i^2=\hat{\bm\epsilon}_i\T\hat{\bm\epsilon}_i/(T-K-1)$,
$\hat\Omega_f$ the sample covariance of the factors and $\hat\Sigma_x=\hat B\hat\Omega_f\hat B\T+\hat\Psi$ \ts{13}{15}{22:10}.
If residual variances differ or residuals are correlated, generalised least squares is the efficient
estimator \ts{13}{15}{23:15}.

\begin{coursenote}[Erratum: slide 11 of \file{lecnote15}]
The estimator of the multifactor covariance is written $\hat\sigma_M^2\hat B\hat\Omega_f\hat B'+\hat\Psi$; the factor
$\hat\sigma_M^2$ is a leftover of the single-index slide. Correct: $\hat\Sigma_x=\hat B\hat\Omega_f\hat B\T+\hat\Psi$ (the factor
variances are already inside $\hat\Omega_f$).
\end{coursenote}

\subsection{Fundamental factor models}
Here the \emph{loadings} are observed asset attributes (industry, market capitalisation, dividend yield,
book-to-price, earnings-to-price, \dots) and the factor realisations are estimated. This is the BARRA approach
of Barr Rosenberg \ts{13}{15}{24:17}, still the standard of commercial risk models.

\paragraph{Industry factor model} \ts{13}{15}{29:35}. Split the $m$ stocks into $K$ industries and let
$\beta_{i,k}=1$ if stock $i$ is in industry $k$, $0$ otherwise: $B$ is a known, time-invariant matrix of
indicators \ts{13}{15}{30:37}. Oil stocks move together, utilities (regulated, low volatility) move together
and differently. For each $t$, estimate $\bm f_t$ by the cross-sectional regression of $\bm x_t$ on $B$
\ts{13}{15}{31:40}:
\begin{equation}\label{ch07:eq-fhat}
  \hat{\bm f}_t=(B\T B)^{-1}B\T\bm x_t,\qquad B\T B=\diag(m_1,\dots,m_K),
\end{equation}
where $m_k$ is the number of stocks in industry $k$: $\hat f_{k,t}$ is simply the \emph{average return of industry $k$}
\ts{13}{15}{32:43}. Then $\hat{\bm\epsilon}_t=\bm x_t-B\hat{\bm f}_t$, $\hat\Omega_f$ and $\hat\Psi$ are sample (co)variances of the
$\hat{\bm f}_t$ and of the residuals, and $\hat\Sigma=B\hat\Omega_fB\T+\hat\Psi$ \ts{13}{15}{33:48}.

\paragraph{Heteroscedasticity.} Stocks have different specific variances, so OLS in \eqref{ch07:eq-fhat} is
inefficient; the GLS estimator weights each stock by the inverse of its noise,
\[
  \hat{\bm f}_t=(B\T\Psi^{-1}B)^{-1}B\T\Psi^{-1}\bm x_t ,
\]
with $\Psi$ estimated in a first pass. Kempthorne stresses the analogy \ts{13}{15}{34:53}: mean-variance
optimisation weights assets by $\Sigma^{-1}\bm\mu$ (reward penalised by variance), GLS weights observations by the
inverse noise variance: in both cases noisy inputs must count less.

\begin{coursenote}[Erratum: slide 17 of \file{lecnote15}]
``$\hat\Sigma=B'\hat\Omega_fB+\hat\Psi$'' has the transposes swapped ($B'\hat\Omega_fB$ is not even defined, $B$ being $m\times K$);
correct: $\hat\Sigma=B\hat\Omega_fB\T+\hat\Psi$.
\end{coursenote}

\paragraph{Factor-mimicking portfolios} \ts{13}{15}{35:56}. The estimate $\hat{\bm f}_t=H\bm x_t$ is linear in the returns
\ts{13}{15}{37:00}, so each row of $H$ is a portfolio of the $m$ assets. Normalising the row weights to sum to one
(one unit of capital) gives a tradeable portfolio whose return tracks the factor \ts{13}{15}{39:04}.

\begin{proposition}[Mimicking portfolios have unit exposure]\label{ch07:prop-mimic}
Let $H=(B\T\Psi^{-1}B)^{-1}B\T\Psi^{-1}$ and let $\bm h_k\T$ be its $k$-th row. Then $HB=I_K$, so under \eqref{ch07:eq-lfm}
with $\bm\alpha=\bm0$ the portfolio $\bm h_k$ has return $f_{k,t}+\bm h_k\T\bm\epsilon_t$: unit exposure to factor $k$ and zero to
the others. Among all portfolios with these exposures ($B\T\bm h=\bm e_k$) it has the smallest specific variance
$\bm h\T\Psi\bm h$ (not discussed in class).
\end{proposition}
\begin{proof}
$HB=(B\T\Psi^{-1}B)^{-1}B\T\Psi^{-1}B=I_K$. For the minimum, the Lagrangian $\frac12\bm h\T\Psi\bm h-\bm\ell\T(B\T\bm h-\bm e_k)$ gives
$\bm h=\Psi^{-1}B\bm\ell$ and $B\T\Psi^{-1}B\bm\ell=\bm e_k$, i.e.\ $\bm h=\Psi^{-1}B(B\T\Psi^{-1}B)^{-1}\bm e_k=\bm h_k$ (Gauss--Markov in
portfolio language).
\end{proof}
In the industry model with OLS, $\bm h_k$ is the equal-weighted portfolio of industry $k$.

\paragraph{Fama--French factors} \ts{13}{15}{25:23}. Fama and French turned attributes into factor
\emph{returns} by sorting. At each $t$, rank the stocks by an attribute (market cap for \emph{size},
book-to-market for \emph{value versus growth}) \ts{13}{15}{26:25}, split into quintiles, and define the factor
as the return of the equal-weighted top quintile minus the bottom quintile (a zero-cost hedge portfolio)
\ts{13}{15}{27:32}. Then estimate each stock's loadings on these factors by time-series regressions. In the
published factors the sign is chosen so that the factor has a positive average return, e.g.\ small minus big
(SMB) and high minus low book-to-market (HML) \ts{13}{15}{28:35}. Value stocks are cheap relative to book
equity; growth stocks are expensive because the price anticipates future growth.

% ------------------------------------------------------------------
\section{Statistical factor models}\label{ch07:sec-stat}
Now neither factors nor loadings are observed: only the returns $\bm x_1,\dots,\bm x_T$ \ts{13}{15}{40:18}. Both
methods below model the sample covariance matrix.

\subsection{Factor analysis by maximum likelihood}\label{ch07:sec-fa}
Take the normalised model with Gaussian variables \ts{13}{15}{47:02}:
$\bm f_t\sim N_K(\bm0,I_K)$, $\bm\epsilon_t\sim N_m(\bm0,\Psi)$, independent over time, so that
$\bm x_t\sim N_m(\bm\alpha,\Sigma_x)$ i.i.d.\ with $\Sigma_x=BB\T+\Psi$ \ts{13}{15}{48:05}. The log-likelihood is
\ts{13}{15}{49:05}
\begin{equation}\label{ch07:eq-loglik}
  \ell(\bm\alpha,\Sigma_x)=-\frac{Tm}{2}\ln(2\pi)-\frac T2\ln\det\Sigma_x-\frac12\sum_{t=1}^T(\bm x_t-\bm\alpha)\T\Sigma_x^{-1}(\bm x_t-\bm\alpha),
\end{equation}
to be maximised over $\bm\alpha,B,\Psi$ subject to $\Sigma_x=BB\T+\Psi$. The maximiser of $\bm\alpha$ is $\bar{\bm x}$, and then
$\ell=-\frac T2\bigl[m\ln2\pi+\ln\det\Sigma_x+\tr(\Sigma_x^{-1}S)\bigr]$ with $S$ the (divisor-$T$) sample covariance.

\begin{coursenote}[Erratum: slide 24 of \file{lecnote15}]
The log-likelihood is written with $-\frac{TK}{2}\log(2\pi)-\frac K2\log|\Sigma|$; taking the logarithm of the likelihood on
slide 23 gives $-\frac{Tm}{2}\ln(2\pi)-\frac T2\ln|\Sigma_x|$, as in \eqref{ch07:eq-loglik}.
\end{coursenote}

\paragraph{The EM algorithm} \ts{13}{15}{51:06}. The likelihood is complicated because the factors are hidden;
if they were observed, estimating $B$ and $\Psi$ would be ordinary regression. The EM algorithm (Dempster, Laird
and Rubin 1977) alternates the two: estimate the hidden variables given the parameters (E-step), re-estimate the
parameters as if the hidden variables were known (M-step); each iteration increases the likelihood
\ts{13}{15}{52:10}. For centred data (not discussed in class; I verified numerically that the likelihood increases
monotonically):
\begin{align*}
  \text{E: }\ & \bm\beta=B\T\Sigma_x^{-1},\quad \E[\bm f_t\mid\bm x_t]=\bm\beta\bm x_t,\quad
     \E[\bm f_t\bm f_t\T\mid\bm x_t]=I-\bm\beta B+\bm\beta\bm x_t\bm x_t\T\bm\beta\T;\\
  \text{M: }\ & B_{\rm new}=S\bm\beta\T\bigl(I-\bm\beta B+\bm\beta S\bm\beta\T\bigr)^{-1},\quad
     \Psi_{\rm new}=\diag\bigl(S-B_{\rm new}\bm\beta S\bigr).
\end{align*}
Time dependence can be added through the linear state-space framework of the Kalman filter
(\cref{ch:timeseries}).

\paragraph{Factor scores} \ts{13}{15}{53:16}. With $\hat{\bm\alpha},\hat B,\hat\Psi$ in hand, estimate each $\bm f_t$ by the
cross-sectional GLS regression of $\bm x_t-\hat{\bm\alpha}$ on $\hat B$:
$\hat{\bm f}_t=(\hat B\T\hat\Psi^{-1}\hat B)^{-1}\hat B\T\hat\Psi^{-1}(\bm x_t-\hat{\bm\alpha})$. These estimated factor series can feed other
models, e.g.\ statistical factors of commodity returns used to explain stocks \ts{13}{15}{54:17}, and can be
rescaled into factor-mimicking portfolios.

\paragraph{Rotations and interpretation} \ts{13}{15}{56:23}. Factor analysis was born in psychometrics:
Spearman's analysis of intelligence tests, where one looks for the ``dimensions of intelligence''. The
orthogonal indeterminacy $B\mapsto BQ$ is used to choose rotated loadings that are easier to interpret (e.g.\ the
varimax criterion, which pushes each loading towards $0$ or $\pm$ large) \ts{13}{15}{57:30}.

\paragraph{How many factors? The likelihood-ratio test} \ts{13}{15}{58:36}. Compare the unrestricted MLE
($\tilde{\bm\alpha}=\bar{\bm x}$, $\tilde\Sigma=S$) with the $K$-factor MLE:
\[\begin{gathered}
  \mathrm{LR}(K)=2\bigl[\ell(\tilde{\bm\alpha},\tilde\Sigma)-\ell(\hat{\bm\alpha},\hat B\hat B\T+\hat\Psi)\bigr]\ \approx\ \chi^2_\nu,\\
  \nu=\frac{m(m+1)}2-\Bigl[mK+m-\frac{K(K-1)}2\Bigr]=\frac{(m-K)^2-(m+K)}2 ,
\end{gathered}\]
where the $K$-factor model has $mK+m$ parameters minus $K(K-1)/2$ for the rotational freedom. $H_0$ ($K$ factors
suffice) is rejected for large $\mathrm{LR}$. For $m=9$: $\nu=12,6,1$ for $K=3,4,5$, and $\nu<0$ for $K=6$, so at most five
factors can be fitted (R's \texttt{factanal} also applies Bartlett's small-sample correction). Unlike PCA, the ML
solution is scale-equivariant: analysing covariances or correlations gives equivalent answers (not discussed in
class).

\subsection{PCA as an approximate factor model; the principal factor method}\label{ch07:sec-pcafa}
Partition the PCA of $\Sigma_x$ into the first $K$ and the remaining $m-K$ components \ts{13}{15}{1:03:46}:
\[
  \bm x=\bm\alpha+\Gamma_1\bm p_1+\Gamma_2\bm p_2=\bm\alpha+B\bm f+\bm\epsilon,\qquad B=\Gamma_1,\ \bm f=\bm p_1,\ \bm\epsilon=\Gamma_2\bm p_2 .
\]
This looks like a factor model (and $B=\Gamma_1\Lambda_1^{1/2}$, $\bm f=\Lambda_1^{-1/2}\bm p_1$ gives orthonormal factors), but
$\Cov\bm\epsilon=\Gamma_2\Lambda_2\Gamma_2\T$ is in general \emph{not diagonal} \ts{13}{15}{1:04:59}: the ``specific'' parts are
correlated across assets. PCA is an approximate factor model, ``principal components'', not ``principal
factors''. The \emph{principal factor method} fixes this by iterating PCA on a reduced covariance:
\begin{enumerate}[1.]
  \item PCA of the sample covariance, $\hat\Sigma_x=\hat\Gamma\hat\Lambda\hat\Gamma\T$;
  \item initial estimates $\tilde B_0=\hat\Gamma_{(K)}\hat\Lambda_{(K)}^{1/2}$ (first $K$ columns and the leading $K\times K$ block),
        $\tilde\Psi_0=\diag(\hat\Sigma_x)-\diag(\tilde B_0\tilde B_0\T)$;
  \item eigen-decompose the reduced matrix $\hat\Sigma_x-\tilde\Psi_0$ (which under the model equals $BB\T$, of rank $K$) and
        repeat step 2 with its eigenvectors and eigenvalues, obtaining $\tilde B_1,\tilde\Psi_1$;
  \item iterate until $\tilde\Psi_s$ stops changing; use the last estimates.
\end{enumerate}

\begin{center}\small
\begin{tabularx}{\linewidth}{@{}l X X@{}}
\toprule
 & \textbf{PCA} & \textbf{Factor analysis (ML)}\\
\midrule
Model & none: exact rotation of the data & $\Sigma_x=BB\T+\Psi$, $\Psi$ diagonal, Gaussian likelihood\\
Output & ordered orthogonal directions, variances & loadings up to rotation, specific variances\\
Residuals & correlated across variables & uncorrelated by assumption\\
Scale & covariance and correlation PCA differ & scale-equivariant\\
Number of factors & scree plot, variance explained & likelihood-ratio test\\
Computation & one SVD & iterative (EM)\\
\bottomrule
\end{tabularx}
\end{center}

\begin{coursenote}[Erratum: slide 32 of \file{lecnote15}]
For the SVD $X^*=VDU'$ of the $m\times T$ centred matrix, $U$ must be $T\times m$ with orthonormal columns, $U'U=I_m$
(the slide says ``$U$: $(m\times T)$ row-orthonormal, $UV'=I_m$'', which is dimensionally inconsistent with
$X^*=VDU'$). With this correction the exercise on the slide (\cref{ch07:ex-lecsvd}) works out.
\end{coursenote}

\begin{casestudy}[2013 Case Study 7: factor modelling of Treasury yields]
\textbf{Data.} FRED constant-maturity yields, 9 tenors (3m to 20y), 3 Jan 2000 to 31 May 2013; daily changes (in
percentage points) over two periods \ts{13}{15}{1:11:14}: 2001--2005 and 2009--May 2013 (the section title says
``2009--2012'', the code uses 2009--2013).
\textbf{2001--2005} \ts{13}{15}{1:12:18}: all mean changes negative (rates fell on balance) \ts{13}{15}{1:13:23}; daily
volatilities 3.8 bp (3m) to 7.0 bp (3y); correlations decreasing away from the diagonal. PCA (\texttt{princomp}):
84.9\%, 9.2\%, 2.9\% \ts{13}{15}{1:18:43}; PC1 a weighted average peaking at 3--5y, PC2 long minus short, PC3
curvature, PC9 the 7y-versus-5y/10y hedge; the PC scores are exactly uncorrelated. Factor analysis
(\texttt{factanal}): 4 factors rejected ($\chi^2=26.8$, 6 df, $p=0.0002$), 5 factors rejected at 5\% ($\chi^2=4.33$, 1 df,
$p=0.037$), 6 factors cannot be fitted ($\nu<0$): no reduced-dimension model passes the test.
\textbf{2009--2013}: short rates at zero, so the 3m--1y yields hardly move (3m volatility 1.25 bp versus 6.6 bp at
7--20y); PCA 88.6\%, 6.8\%, 1.6\%, with PC1 loadings near zero at the short end; PC2 now opposes long and
\emph{medium} tenors. Factor analysis: 3 factors rejected ($\chi^2=231$, 12 df), 4 factors accepted ($\chi^2=5.2$, 6 df,
$p=0.52$); the 3m yield has a large specific variance (uniqueness 0.65).
\textbf{Lesson} \ts{13}{15}{1:25:03}: the structure depends on the period (monetary regime), so the calibration
window matters, and the fitted factors are only a starting point for modelling their dynamics (as in 2013 PS7
Problem~4). (Several uniquenesses equal 0.005, the lower bound used by \texttt{factanal}: near-``Heywood'' cases,
where a variable is almost entirely common factor; not discussed in class.)
\end{casestudy}

\begin{coursenote}[2013 versus 2024]
Both years analyse daily changes of US Treasury yields and find level, slope and curvature. 2013 (Kempthorne,
statistics) uses FRED constant-maturity yields at 9 tenors, fixed 5-year windows, and adds the factor-model and
factor-analysis machinery; 2024 (Andreev, practice) uses the Fed's fitted par yields at 30 maturities with a
rolling 2-year window, and focuses on stability, out-of-sample behaviour and PC-neutral trading. The 2024 course
has no lecture on macroeconomic or fundamental factor models; their portfolio-theory side reappears in
\cref{ch:portfolio}.
\end{coursenote}

% ------------------------------------------------------------------
\begin{keyformulas}
\begin{align*}
&\text{PCA:} && \Sigma=\Gamma\Lambda\Gamma\T,\quad
  \bm p=\Gamma\T(\bm x-\bm\mu),\quad \Cov\bm p=\Lambda,\\
& && \bm x=\bm\mu+\textstyle\sum_kp_k\bm\gamma_k\\
&\text{Variational:} && \bm\gamma_1=\arg\max_{\|\bm w\|=1}\bm w\T\Sigma\bm w,\\
& && \text{share}_k=\lambda_k/\tr\Sigma\\
&\text{Best subspace:} && \min_V\E[\|(I-VV\T)(\bm x-\bm\mu)\|^2]\\
& && =\textstyle\sum_{j>K}\lambda_j\ \ (V=\Gamma_K)\\
&\text{Sample PCA:} && X_c=UD\hat\Gamma\T,\quad
  \hat\lambda_k=d_k^2/(N-1),\quad P=X_c\hat\Gamma=UD\\
&\text{Correlation PCA:} && R=D_\sigma^{-1}\Sigma D_\sigma^{-1},\quad \tr R=p\\
&\text{Out of sample:} && W\T\Sigma^YW=D\T\Lambda^YD,\\
& && D=(W^Y)\T W;\quad \mathrm{ED}=(\textstyle\sum\lambda)^2/\textstyle\sum\lambda^2\\
&\text{PC exposures:} && e_k=\bm h\T\bm\gamma_k;\\
& && \text{curve } \alpha=\gamma_{a1}/\gamma_{b1};\\
& && \text{fly: } 2\times2\text{ system \eqref{ch07:eq-fly}}\\
&\text{PC1 hedge:} && \bm V'=\bm V-(\bm V\T\bm\gamma_1)\bm\gamma_1,\\
& && \Var=\bm V\T\Sigma\bm V-\lambda_1(\bm V\T\bm\gamma_1)^2\\
&\text{Factor model:} && \bm x_t=\bm\alpha+B\bm f_t+\bm\epsilon_t,\\
& && \Sigma_x=B\Omega_fB\T+\Psi\ \ (\Psi\text{ diagonal})\\
&\text{Single index:} && \Sigma_x=\sigma_M^2\bm\beta\bm\beta\T+\Psi\\
&\text{Factor estimates:} && \hat{\bm f}_t=(B\T\Psi^{-1}B)^{-1}B\T\Psi^{-1}\bm x_t,\\
& && HB=I_K\\
&\text{LR test df:} && \nu=\tfrac12\bigl[(m-K)^2-(m+K)\bigr]
\end{align*}
\end{keyformulas}

\begin{exercise}[not from the course]\label{ch07:ex-fa}
\leavevmode
\begin{enumerate}[(a)]
  \item Count the free parameters of the $K$-factor model $\Sigma_x=BB\T+\Psi$ and derive the degrees of freedom $\nu$ of the
      likelihood-ratio test.
  \item Check the values of $\nu$ reported in 2013 Case Study~7 and explain why a 6-factor model
      cannot be fitted to 9 yields.
  \item For the single-index model with $\Psi=\psi I$, show that PCA recovers $\bm\beta/\|\bm\beta\|$ as
      PC1; what happens if the $\sigma_i^2$ differ?
\end{enumerate}
\end{exercise}

% ------------------------------------------------------------------
\section*{Exercises}
\addcontentsline{toc}{section}{Exercises}
Standalone exercises follow the section they practise. The exercises that build on one another stay together here: PS3 Problem~4 with the exercises that reuse its data (PS7 Problem~4 and the PC-neutral hedges), and PS7 Problems~1--2.

\subsection*{From 2024 Problem Set 3}

\begin{exercise}[PS3 (2024), Problem 4; also PS7 (2013), Problem 3: PCA of yield changes]\label{ch07:ex-ps3}
Daily changes of the 9 FRED constant-maturity Treasury yields over 2001--2005 are analysed twice, with the sample
covariance matrix and with the sample correlation matrix. The standard deviations of the PC variables are
\begin{center}\small
\begin{tabular}{@{}l rrrrrrrrr@{}}
\toprule
 & 1 & 2 & 3 & 4 & 5 & 6 & 7 & 8 & 9\\
\midrule
covariance & 0.1618 & 0.0532 & 0.0299 & 0.0192 & 0.0135 & 0.0112 & 0.0096 & 0.0090 & 0.0081\\
correlation & 2.6338 & 1.1892 & 0.5652 & 0.3971 & 0.2694 & 0.2080 & 0.1485 & 0.1366 & 0.1225\\
\bottomrule
\end{tabular}
\end{center}
(proportions of variance 84.9, 9.2, 2.9\,\% and 77.1, 15.7, 3.5\,\% for the first three) and the first three loading
vectors are
\begin{center}\small
\begin{tabular}{@{}l rrr c rrr@{}}
\toprule
 & \multicolumn{3}{c}{covariance PCA} & & \multicolumn{3}{c}{correlation PCA}\\
\cmidrule(lr){2-4}\cmidrule(l){6-8}
Tenor & PC1 & PC2 & PC3 & & PC1 & PC2 & PC3\\
\midrule
3M  & 0.1088 & 0.5225 & 0.5434 & & $-0.2083$ & 0.6293 & 0.5831\\
6M  & 0.1581 & 0.4817 & 0.2508 & & $-0.2830$ & 0.5154 & $-0.0654$\\
1Y  & 0.2531 & 0.4107 & $-0.0514$ & & $-0.3369$ & 0.2698 & $-0.4066$\\
2Y  & 0.3905 & 0.2071 & $-0.4750$ & & $-0.3625$ & 0.0064 & $-0.4005$\\
3Y  & 0.4197 & 0.0873 & $-0.4022$ & & $-0.3664$ & $-0.0665$ & $-0.2719$\\
5Y  & 0.4206 & $-0.1082$ & $-0.0311$ & & $-0.3670$ & $-0.1619$ & 0.0036\\
7Y  & 0.4009 & $-0.2289$ & 0.1475 & & $-0.3608$ & $-0.2298$ & 0.1435\\
10Y & 0.3697 & $-0.2837$ & 0.2710 & & $-0.3525$ & $-0.2656$ & 0.2554\\
20Y & 0.3101 & $-0.3621$ & 0.3944 & & $-0.3290$ & $-0.3338$ & 0.4125\\
\bottomrule
\end{tabular}
\end{center}
(a) PC1 loadings are all positive in the covariance case and all negative in the correlation case. Is the
difference in sign meaningful? (Does the eigen-decomposition change if an eigenvector is multiplied by $-1$?)
(b) The PC1 loadings range more widely in the covariance case; the loading on the least variable change (3M) is
larger in magnitude in the correlation case, and those on the most variable changes are smaller. Explain why,
recalling that PC1 is the unit-norm weighted combination with the largest variance and that in the correlation
case each yield change is scaled to unit variance.
(c) Interpret the first three PCs of the correlation case and compare with the covariance case.
(d) If the goal is to model the dynamics of the whole term structure across all tenors, argue why correlation PCA
might be preferred.
(e) (2024 only) Compute the total variance of the correlation-case PCs. Is it an integer? Explain.
\end{exercise}

\begin{coursenote}[On the data of the previous exercise]
The problem sets say the changes are ``in basis point units''; the numbers are in percentage points (a PC1
standard deviation of 0.1618 is 16 bp per day, the tenor volatilities of Case Study 7 are 0.04--0.07). The 2013
version prints correlation-case standard deviations about $0.04\%$ smaller (their squares sum to $8.993$ instead of
$9$, a factor $(N-1)/N$ with $N\approx1247$ days, presumably a different divisor convention); the loadings coincide up
to signs.
\end{coursenote}

\subsection*{From 2013 Problem Set 7}

\begin{exercise}[PS7 (2013), Problem 1: PCA of a bivariate vector]\label{ch07:ex-ps7-1}
Let $\bm X=(X_1,X_2)\T$ have mean $\bm\alpha=(\alpha_1,\alpha_2)\T$ and covariance
$\Sigma=\bigl(\begin{smallmatrix}\sigma_1^2&\rho\sigma_1\sigma_2\\ \rho\sigma_1\sigma_2&\sigma_2^2\end{smallmatrix}\bigr)$.
\begin{enumerate}[(a)]
  \item Compute the eigenvalues $\lambda_1\ge\lambda_2\ge0$ of $\Sigma$.
  \item Compute orthonormal eigenvectors $\bm\gamma_1,\bm\gamma_2$.
  \item Verify $\Sigma=\sum_i\lambda_i\bm\gamma_i\bm\gamma_i\T$.
  \item For $p_i=\bm\gamma_i\T(\bm X-\bm\alpha)$ prove $\E[p_i]=0$, $\Var(p_i)=\lambda_i$ and $\Cov(p_1,p_2)=0$.
  \item Are the eigenvectors in (b) unique? If not, how are different solutions related?
\end{enumerate}
\end{exercise}

\begin{exercise}[PS7 (2013), Problem 2: PCA as a distance-preserving affine map]\label{ch07:ex-ps7-2}
A distance-preserving affine map of the plane can be written $T(\bm X)=\Phi\T(\bm X-\bm\mu)$ with $\Phi=[\bm\phi_1:\bm\phi_2]$
orthogonal and $\bm\mu\in\R^2$: it moves the origin to $\bm\mu$ and rotates the axes.
\begin{enumerate}[(a)]
  \item Find the unit vector $\bm\phi_1$ maximising $\Var(\bm\phi_1\T\bm X)$ and show that the maximum is $\lambda_1$ of the previous
      exercise.
  \item Find the unit vector $\bm\phi_2$ minimising $\Var(\bm\phi_2\T\bm X)$ and show that the minimum is $\lambda_2$.
  \item Give an equivalent definition of PCA in terms of affine transformations of the variables and their
      variances and expectations.
\end{enumerate}
\end{exercise}

\begin{exercise}[PS7 (2013), Problem 4: dynamics of the PC scores]\label{ch07:ex-ps7-4}
A VAR(3) with constant (\cref{ch:timeseries}) was fitted to the scores of the first three PCs of the
correlation-matrix analysis of \cref{ch07:ex-ps3} ($T=1245$ days). Estimates and $t$-values:
\begin{center}\small
\begin{tabular}{@{}l rr c rr c rr@{}}
\toprule
 & \multicolumn{2}{c}{eq.\ Comp.1} & & \multicolumn{2}{c}{eq.\ Comp.2} & & \multicolumn{2}{c}{eq.\ Comp.3}\\
\cmidrule(lr){2-3}\cmidrule(lr){5-6}\cmidrule(l){8-9}
Regressor & est. & $t$ & & est. & $t$ & & est. & $t$\\
\midrule
Comp.1, lag 1 & 0.0527 & 1.86 & & $-0.0203$ & $-1.65$ & & $-0.0029$ & $-0.48$\\
Comp.2, lag 1 & $-0.1340$ & $-2.08$ & & 0.1382 & 4.93 & & $-0.0143$ & $-1.03$\\
Comp.3, lag 1 & $-0.1322$ & $-1.00$ & & $-0.0833$ & $-1.45$ & & 0.0163 & 0.58\\
Comp.1, lag 2 & $-0.0530$ & $-1.87$ & & $-0.0072$ & $-0.59$ & & $-0.0135$ & $-2.23$\\
Comp.2, lag 2 & $-0.0542$ & $-0.84$ & & 0.0370 & 1.32 & & $-0.0037$ & $-0.27$\\
Comp.3, lag 2 & 0.1825 & 1.39 & & $-0.0195$ & $-0.34$ & & $-0.0523$ & $-1.86$\\
Comp.1, lag 3 & $-0.0169$ & $-0.60$ & & $-0.0085$ & $-0.69$ & & $-0.0019$ & $-0.32$\\
Comp.2, lag 3 & 0.0130 & 0.21 & & $-0.0403$ & $-1.46$ & & $-0.0030$ & $-0.22$\\
Comp.3, lag 3 & 0.0830 & 0.63 & & 0.0148 & 0.26 & & $-0.0712$ & $-2.53$\\
constant & $-0.0114$ & $-0.15$ & & 0.0121 & 0.38 & & 0.0002 & 0.01\\
\midrule
$R^2$; $F$-test $p$-value & 0.013 & 0.054 & & 0.029 & $4\cdot10^{-5}$ & & 0.013 & 0.059\\
\bottomrule
\end{tabular}
\end{center}
Residual standard errors are 2.60, 1.13, 0.56; residual correlations are below $0.025$ in absolute value; the
largest root of the characteristic polynomial has modulus 0.438.
(a)--(c) Interpret the estimated equations for Comp.1, Comp.2 and Comp.3.
(d) For each equation, is there evidence of mean reversion or of momentum in the PC variable? Explain.
\end{exercise}

\subsection*{Additional exercises}

\begin{exercise}[not from the course]\label{ch07:ex-fly}
Using the covariance-PCA loadings of \cref{ch07:ex-ps3}:
\begin{enumerate}[(a)]
  \item find the PC1-neutral 2s10s steepener and check
      that its exposure to PC2 is non-zero;
  \item find the PC1/PC2-neutral 2s5s10s fly (buy \$1 dv01 of 5y) and compute its
      exposure to PC3;
  \item prove the variance formula for the PC1-hedged portfolio $\bm V'=\bm V-(\bm V\T\bm\gamma_1)\bm\gamma_1$ and show
      that it never exceeds $\Var(\bm V\T\bm r)$.
\end{enumerate}
\end{exercise}
